\documentclass[12pt,a4paper]{article}

%!TEX program = pdflatex

\usepackage[utf8]{inputenc}
\usepackage[T5]{fontenc}
\usepackage[vietnamese]{babel}
\usepackage[margin=2.5cm]{geometry}
\usepackage{setspace}
\onehalfspacing
\usepackage{graphicx}
\usepackage{tikz}
\usetikzlibrary{arrows.meta, positioning, shapes.geometric}
\usepackage{array}
\usepackage{booktabs}
\usepackage{multirow}
\usepackage{longtable}
\usepackage{float}
\usepackage{enumitem}
\usepackage{tabularx}
\usepackage{makecell}
\usepackage{ragged2e}
\usepackage{xcolor}
\usepackage{titlesec}
\usepackage{tocloft}
\usepackage{xurl}
\usepackage{seqsplit}
\usepackage{hyperref}

\hypersetup{
    colorlinks=false,
    hidelinks,
    pdftitle={Đặc tả và phân tích yêu cầu chức năng phần mềm - Hệ thống giúp giáo viên quản lý lớp học lập trình thi đấu},
    pdfauthor={Nhóm phát triển hệ thống}
}

\setlength{\parindent}{0pt}
\setlength{\parskip}{6pt}
\setlength{\emergencystretch}{3em}
\setcounter{secnumdepth}{3}
\raggedbottom
\sloppy

\renewcommand{\thesection}{\Roman{section}}
\renewcommand{\thesubsection}{\arabic{subsection}}
\renewcommand{\thesubsubsection}{\thesubsection.\arabic{subsubsection}}

\titleformat{\section}{\large\bfseries}{\thesection.}{0.5em}{}
\titleformat{\subsection}{\normalsize\bfseries}{\thesubsection.}{0.5em}{}
\titleformat{\subsubsection}{\normalsize\itshape\bfseries}{\thesubsubsection.}{0.5em}{}
\titleformat{\paragraph}[block]{\normalsize\bfseries}{}{0pt}{}
\titlespacing*{\paragraph}{0pt}{6pt}{0pt}

\renewcommand{\cftsecaftersnum}{.}
\renewcommand{\cftsubsecaftersnum}{.}
\renewcommand{\cftsubsubsecaftersnum}{.}

\definecolor{uitblue}{HTML}{005BAA}
\definecolor{uitlightblue}{HTML}{EAF4FF}

\tikzset{
    dfdprocess/.style={ellipse, draw=uitblue, line width=0.8pt, minimum width=3.6cm, minimum height=1.35cm, align=center, fill=uitlightblue},
    dfdentity/.style={rectangle, draw=black!70, rounded corners=2pt, minimum width=2.8cm, minimum height=0.9cm, align=center, fill=white},
    dfddatastore/.style={rectangle, draw=black!70, minimum width=3.5cm, minimum height=0.75cm, align=center, fill=white},
    dfdarrow/.style={-Stealth, semithick},
    dfdbluearrow/.style={-Stealth, semithick, draw=uitblue}
}
\newcommand{\dfddiagram}[4]{%
\begin{center}
\resizebox{0.96\textwidth}{!}{%
\begin{tikzpicture}[node distance=1.8cm and 2.6cm]
    \node[dfdprocess] (process) {#1};
    \node[dfdentity, above=of process] (actor) {#2};
    \node[dfdentity, left=of process] (input) {#3};
    \node[dfdentity, right=of process] (output) {Thiết bị xuất};
    \node[dfddatastore, below=of process] (storage) {#4};

    \draw[dfdarrow] ([xshift=-0.4cm]actor.south) -- node[left] {D1} ([xshift=-0.4cm]process.north);
    \draw[dfdbluearrow] ([xshift=0.4cm]process.north) -- node[right] {D6} ([xshift=0.4cm]actor.south);
    \draw[dfdbluearrow] (input.east) -- node[above] {D2} (process.west);
    \draw[dfdarrow] (process.east) -- node[above] {D5} (output.west);
    \draw[dfdarrow] ([xshift=-0.4cm]storage.north) -- node[left] {D3} ([xshift=-0.4cm]process.south);
    \draw[dfdarrow] ([xshift=0.4cm]process.south) -- node[right] {D4} ([xshift=0.4cm]storage.north);
\end{tikzpicture}%
}
\end{center}
}

\newcolumntype{C}[1]{>{\Centering\hspace{0pt}\arraybackslash}p{#1}}
\newcolumntype{L}[1]{>{\RaggedRight\hspace{0pt}\arraybackslash}p{#1}}
\newcolumntype{Y}{>{\RaggedRight\hspace{0pt}\arraybackslash}X}
\newcommand{\tablesetup}{\scriptsize\setlength{\tabcolsep}{2pt}\renewcommand{\arraystretch}{1.2}}
\newcommand{\code}[1]{{\ttfamily\seqsplit{#1}}}
\newcommand{\uitlogo}{%
\IfFileExists{logo_uit.jpg}{%
    
\includegraphics[width=0.28\textwidth]{logo_uit.jpg}%
}{%
\IfFileExists{./logo_uit.jpg}{%
    
\includegraphics[width=0.28\textwidth]{./logo_uit.jpg}%
}{%
\IfFileExists{../logo_uit.jpg}{%
    
\includegraphics[width=0.28\textwidth]{../logo_uit.jpg}%
}{%
\IfFileExists{../../logo_uit.jpg}{%
    
\includegraphics[width=0.28\textwidth]{../../logo_uit.jpg}%
}{%
    \fcolorbox{uitblue}{white}{\parbox[c][2.7cm][c]{0.28\textwidth}{\centering\textcolor{uitblue}{Logo UIT}}}%
}}}}%
}

\begin{document}

\begin{titlepage}
    \centering
    \begin{tikzpicture}[remember picture,overlay]
        \fill[uitlightblue] (current page.south west) rectangle (current page.north east);
        \draw[uitblue,line width=2pt]
            ([xshift=0.65cm,yshift=0.65cm]current page.south west)
            rectangle
            ([xshift=-0.65cm,yshift=-0.65cm]current page.north east);
    \end{tikzpicture}

    \vspace*{0.8cm}
    \begin{minipage}[c][0.86\textheight][c]{0.82\textwidth}
        \centering
        \uitlogo\\[0.8cm]

        {\Large \textbf{\textcolor{uitblue}{ĐẠI HỌC CÔNG NGHỆ THÔNG TIN}}}\\[0.15cm]
        {\large \textbf{ĐẠI HỌC QUỐC GIA TP.HCM}}\\[0.9cm]

        {\large \textbf{SE104 - NHẬP MÔN CÔNG NGHỆ PHẦN MỀM}}\\[0.5cm]

        \rule{0.72\textwidth}{0.7pt}\\[0.6cm]
        {\Large \textbf{HỆ THỐNG GIÚP GIÁO VIÊN}}\\[0.2cm]
        {\Large \textbf{QUẢN LÝ LỚP HỌC LẬP TRÌNH THI ĐẤU}}\\[0.6cm]
        \rule{0.72\textwidth}{0.7pt}\\[1.2cm]

        {\large \textbf{Giáo viên hướng dẫn: ThS. Đỗ Văn Tiến}}\\[1.2cm]

        {\large \textbf{Nhóm sinh viên thực hiện}}\\[0.45cm]
        \begin{tabular}{ll}
            Nguyễn Thành Sơn & 24521532\\
            Lê Quang Trung & 24521883\\
            Đinh Mạnh Hùng & 24520014\\
            Trần Vạn Tấn & 23521407\\
        \end{tabular}

        \vfill
        {\large \textbf{Lớp SE104.Q26}}\\
    \end{minipage}
\end{titlepage}

\tableofcontents
\newpage

\section{Tổng quan đề tài}

\subsection{Giới thiệu nhóm thực hiện}

Nhóm thực hiện đề tài gồm bốn thành viên. Các thành viên cùng tham gia khảo sát
hiện trạng, phân tích yêu cầu, thiết kế nghiệp vụ, xây dựng hệ thống và hoàn
thiện tài liệu báo cáo.

\begin{table}[!htbp]
\centering
\tablesetup
\begin{tabularx}{\textwidth}{|C{0.8cm}|L{4cm}|L{2.6cm}|Y|}
\hline
\textbf{STT} & \textbf{Họ và tên} & \textbf{MSSV} & \textbf{Vai trò chính trong đề tài} \\
\hline
1 & Nguyễn Thành Sơn & 24521532 & Phân tích nghiệp vụ, đặc tả yêu cầu và phối hợp xây dựng các chức năng quản lý lớp học. \\
\hline
2 & Lê Quang Trung & 24521883 & Thiết kế và phát triển các chức năng quản lý học sinh, ghi danh, điểm danh và báo cáo. \\
\hline
3 & Đinh Mạnh Hùng & 24520014 & Thiết kế dữ liệu, xử lý tài chính, công nợ và tích hợp Codeforces Gym. \\
\hline
4 & Trần Vạn Tấn & 23521407 & Xây dựng giao diện, tích hợp Discord, kiểm thử và hoàn thiện tài liệu. \\
\hline
\end{tabularx}
\end{table}

\subsection{Giới thiệu đề tài}

Trong quá trình quản lý lớp dạy thêm thuật toán, giáo viên thường phải theo dõi
nhiều nhóm dữ liệu cùng lúc: hồ sơ học sinh, lớp học, lịch học, buổi học, điểm
danh, học phí, công nợ, thông báo Discord và kết quả luyện tập trên Codeforces.
Nếu các dữ liệu này được lưu rời rạc bằng bảng tính, tin nhắn hoặc ghi chú cá
nhân, giáo viên rất khó có một cái nhìn thống nhất về tình hình học tập và tài
chính của từng học sinh.

Những bất cập chính của cách quản lý bán thủ công gồm:

\begin{itemize}[leftmargin=*]
\item \textbf{Dữ liệu học sinh và ghi danh phân tán:} Hồ sơ học sinh, lớp đang học, lịch sử chuyển lớp hoặc rút khỏi lớp thường nằm ở nhiều file khác nhau, khó truy vết theo thời gian.
\item \textbf{Quản lý buổi học và điểm danh phức tạp:} Lớp học có lịch định kỳ nhưng thực tế có thể phát sinh buổi học thủ công, hủy buổi học hoặc thay đổi trạng thái điểm danh. Nếu xử lý bằng tay, giáo viên dễ bỏ sót dữ liệu.
\item \textbf{Đối soát tài chính dễ sai lệch:} Học phí phát sinh theo buổi học và trạng thái điểm danh, trong khi tiền thu hoặc hoàn tiền lại được ghi nhận theo từng giao dịch. Nếu không có hệ thống liên kết giữa khoản phí và giao dịch, việc tính số dư học sinh dễ nhầm lẫn.
\item \textbf{Thiếu tích hợp Discord:} Việc mời học sinh vào server lớp, chọn channel thông báo, gửi thông báo chung hoặc tin nhắn riêng thường phải thao tác thủ công, khó kiểm soát kết quả gửi.
\item \textbf{Thiếu liên kết với Codeforces Gym:} Codeforces có dữ liệu Gym và bảng xếp hạng, nhưng giáo viên cần một nơi gắn Gym với lớp học, đồng bộ danh sách \code{GymProblem} và xem kết quả theo danh sách học sinh đang học.
\end{itemize}

Vì vậy, đề tài ``Hệ thống giúp giáo viên quản lý lớp học lập trình thi đấu'' được xây dựng nhằm tập
trung hóa dữ liệu vận hành lớp học, hỗ trợ giáo viên theo dõi tiến độ học tập,
điểm danh, công nợ và các tích hợp Discord/Codeforces trong một hệ thống web
thống nhất.

\subsection{Mục tiêu phần mềm}

Phần mềm được xây dựng với các mục tiêu chính sau:

\subsubsection{Quản lý dữ liệu học sinh và ghi danh tập trung}

\begin{itemize}[leftmargin=*]
\item Lưu trữ hồ sơ học sinh, số điện thoại, ghi chú, Codeforces handle và Discord identity nếu học sinh đã uỷ quyền.
\item Quản lý vòng đời học sinh theo trạng thái đang học, chờ lưu trữ và đã lưu trữ.
\item Theo dõi lịch sử ghi danh, chuyển lớp, rút học sinh và khôi phục học sinh vào lớp mới.
\end{itemize}

\subsubsection{Quản lý lớp học, lịch học, buổi học và điểm danh}

\begin{itemize}[leftmargin=*]
\item Cho phép giáo viên tạo lớp học, thiết lập học phí mỗi buổi và lịch học định kỳ theo thứ/giờ.
\item Quản lý các buổi học thực tế, bao gồm buổi học được sinh theo lịch hoặc buổi học tạo thủ công.
\item Hỗ trợ cập nhật điểm danh theo từng học sinh trong buổi học với các trạng thái có mặt, vắng có phép và vắng không phép.
\end{itemize}

\subsubsection{Tự động hóa tài chính và công nợ}

\begin{itemize}[leftmargin=*]
\item Ghi nhận khoản phí phát sinh từ buổi học bằng \code{FeeRecord}.
\item Ghi nhận giao dịch thu tiền hoặc hoàn tiền bằng \code{Transaction}.
\item Tính số dư học sinh, tổng quan tài chính và báo cáo thu nhập theo kỳ để hỗ trợ đối soát.
\end{itemize}

\subsubsection{Tích hợp Discord}

\begin{itemize}[leftmargin=*]
\item Cho phép quản trị viên cấu hình bot Discord và OAuth credential dùng chung.
\item Cho phép giáo viên liên kết tài khoản Discord, chọn Discord guild và channel cho lớp học.
\item Hỗ trợ học sinh uỷ quyền Discord OAuth để hệ thống lưu \code{StudentDiscordCredential} và quản lý membership của học sinh trong guild lớp.
\item Hỗ trợ gửi thông báo vào channel lớp hoặc gửi tin nhắn riêng cho học sinh.
\end{itemize}

\subsubsection{Tích hợp Codeforces Gym}

\begin{itemize}[leftmargin=*]
\item Cho phép giáo viên cấu hình Codeforces credential để đồng bộ dữ liệu Gym.
\item Cho phép gắn Codeforces Gym vào lớp học, đồng bộ danh sách \code{GymProblem} và bảng xếp hạng.
\item Hiển thị kết quả theo ma trận học sinh - \code{GymProblem} để giáo viên theo dõi tiến độ luyện tập.
\end{itemize}

\subsection{Phạm vi phần mềm}

Hệ thống tập trung vào các phân hệ chức năng sau:

\begin{itemize}[leftmargin=*]
\item \textbf{Quản lý tài khoản:} Đăng ký, đăng nhập, cập nhật hồ sơ giáo viên, quản lý tài khoản giáo viên và cấu hình credential cần thiết.
\item \textbf{Quản lý học sinh và ghi danh:} Tạo hồ sơ học sinh, cập nhật thông tin, chuyển lớp, rút học sinh, khôi phục và lưu trữ học sinh.
\item \textbf{Quản lý lớp học, lịch học và buổi học:} Tạo lớp, thiết lập lịch định kỳ, tạo buổi học thủ công, xem danh sách buổi học và hủy buổi học khi cần.
\item \textbf{Quản lý điểm danh:} Xem và cập nhật trạng thái điểm danh theo từng buổi học.
\item \textbf{Quản lý tài chính:} Ghi nhận giao dịch, quản lý khoản phí phát sinh, xem số dư học sinh, tổng quan tài chính và báo cáo thu nhập.
\item \textbf{Quản lý Discord:} Cấu hình bot, xác minh Discord của giáo viên, uỷ quyền Discord của học sinh, liên kết guild/channel cho lớp và gửi thông báo.
\item \textbf{Quản lý Codeforces Gym:} Gắn/gỡ Gym khỏi lớp, đồng bộ \code{GymProblem} và xem bảng xếp hạng Gym.
\item \textbf{Dashboard và báo cáo học tập:} Cung cấp dashboard tổng quan và hồ sơ học tập của từng học sinh dựa trên dữ liệu đang có trong hệ thống.
\end{itemize}

Ngoài phạm vi hiện tại, hệ thống không có module giao homework riêng, không có luồng học sinh nộp bài trực tiếp lên hệ thống, không có chức năng giáo viên chấm điểm hoặc nhận xét bài làm. Việc theo dõi tiến độ lập trình chỉ dựa trên dữ liệu Codeforces Gym đã đồng bộ.

\section{Xác định và thu thập yêu cầu}

\subsection{Kế hoạch và phương pháp khảo sát}

\subsubsection{Đối tượng và địa điểm khảo sát}

- Địa điểm khảo sát: Khảo sát được thực hiện thông qua trao đổi trực tiếp và
trao đổi online với giáo viên đang dạy thêm thuật toán, hiện đang vận hành lớp
học bằng các công cụ rời rạc như Google Sheets, Excel, Discord, Codeforces và
các ứng dụng ghi chú.

- Các đối tượng liên quan cần khảo sát:

\begin{itemize}[leftmargin=*]
\item \textbf{Giáo viên dạy thêm thuật toán:} Người sử dụng chính của hệ thống, trực tiếp quản lý học sinh, lớp học, lịch học, buổi học, điểm danh, tài chính, Discord và Codeforces Gym.
\item \textbf{Quản trị viên hệ thống:} Người cấu hình tài khoản giáo viên và bot Discord dùng chung.
\item \textbf{Học sinh:} Đối tượng liên quan đến uỷ quyền Discord OAuth, tham gia guild lớp và được theo dõi kết quả luyện tập trên Codeforces Gym.
\end{itemize}

\subsubsection{Các phương pháp khảo sát}

- \textbf{Phỏng vấn:} Thu thập thông tin chuyên sâu về quy trình vận hành lớp
học, cách ghi nhận điểm danh, cách tính học phí, cách gửi thông báo và các khó
khăn khi theo dõi kết quả Codeforces.

- \textbf{Quan sát thực tế:} Theo dõi cách giáo viên đang tạo lớp, lập lịch,
điểm danh, ghi nhận học phí, gửi thông báo Discord và xem kết quả luyện tập của
học sinh trên Codeforces.

- \textbf{Nghiên cứu tài liệu:} Phân tích các file bảng tính, danh sách học
sinh, lịch học, bảng công nợ, kênh Discord lớp học và dữ liệu Codeforces Gym mà
giáo viên đang sử dụng.

- \textbf{Tổng hợp nghiệp vụ hiện có:} Chuẩn hóa các thao tác rời rạc thành
các nhóm chức năng ổn định để làm cơ sở đặc tả yêu cầu phần mềm.

\subsection{Khảo sát hiện trạng}

\subsubsection{Hiện trạng tổ chức}

Mô hình tổ chức của hoạt động dạy thêm thuật toán tương đối gọn. Giáo viên là
người vận hành chính, chịu trách nhiệm quản lý học sinh, lớp học, lịch học, buổi
học, điểm danh, học phí, công nợ và theo dõi kết quả luyện tập. Quản trị viên
chỉ tham gia ở các công việc cấu hình hệ thống như quản lý tài khoản giáo viên
và cấu hình bot Discord.

Học sinh không sử dụng toàn bộ hệ thống như giáo viên, nhưng có tương tác trực
tiếp trong luồng uỷ quyền Discord OAuth để hệ thống lưu danh tính Discord và
phục vụ việc tham gia guild lớp học. Discord và Codeforces là hai hệ thống ngoài
đóng vai trò cung cấp dữ liệu hoặc kênh tương tác.

\subsubsection{Hiện trạng nghiệp vụ}

\paragraph{1. Quản lý học sinh và ghi danh}
Thông tin học sinh thường được lưu bằng bảng tính hoặc ghi chú rời rạc. Việc
theo dõi học sinh đang học lớp nào, đã chuyển lớp khi nào, đã rút khỏi lớp hay
đang chờ xử lý công nợ chưa có một quy trình tập trung.

\paragraph{2. Quản lý lớp học, lịch học và buổi học}
Lịch học định kỳ thường được ghi trong bảng tính hoặc tin nhắn ghim. Khi phát
sinh buổi học thủ công, hủy buổi học hoặc thay đổi lịch, giáo viên phải cập
nhật nhiều nơi nên dễ thiếu lịch sử chính xác.

\paragraph{3. Quản lý điểm danh}
Điểm danh thường được ghi bằng tay hoặc dựa vào việc học sinh có mặt trong lớp
online. Dữ liệu điểm danh chưa được liên kết chặt với khoản phí phát sinh, gây
khó khăn khi cần kiểm tra vì sao học sinh phát sinh học phí ở một buổi cụ thể.

\paragraph{4. Quản lý tài chính và công nợ}
Giáo viên phải tự đối chiếu giữa số buổi học, học phí mỗi buổi, khoản phí phải
thu, các giao dịch đã thu và các khoản hoàn tiền nếu có. Quy trình này dễ sai
khi học sinh chuyển lớp, rút khỏi lớp hoặc có nhiều giao dịch nhỏ.

\paragraph{5. Quản lý Discord}
Thông báo lịch học, nhắc học phí hoặc trao đổi với học sinh thường được thực
hiện qua Discord bằng thao tác thủ công. Giáo viên cần một cách chọn guild,
channel thông báo và gửi tin nhắn từ hệ thống để giảm sai sót.

\paragraph{6. Theo dõi Codeforces Gym}
Codeforces cung cấp Gym, \code{GymProblem} và standing, nhưng giáo viên vẫn phải tự mở
Codeforces để đối chiếu. Chưa có nơi gắn Gym với lớp học và xem bảng kết quả
theo đúng danh sách học sinh đang ghi danh.

\paragraph{7. Báo cáo và thống kê}
Báo cáo thu nhập, số dư học sinh, tổng quan lớp học và hồ sơ học tập thường
được tổng hợp thủ công từ nhiều nguồn, mất thời gian và khó kiểm tra lại.

\subsubsection{Hiện trạng tin học}

\paragraph{1. Phần cứng đang sử dụng}
Giáo viên chủ yếu sử dụng laptop cá nhân hoặc máy tính để bàn có kết nối
Internet. Học sinh sử dụng máy tính cá nhân để học, dùng Discord và Codeforces.

\paragraph{2. Phần mềm và công cụ đang sử dụng}
\begin{itemize}[leftmargin=*]
\item Google Sheets hoặc Microsoft Excel để lưu danh sách học sinh, lịch học và công nợ.
\item Discord để quản lý server lớp học, trao đổi và gửi thông báo.
\item Codeforces để tổ chức hoặc theo dõi Gym, \code{GymProblem} và standing.
\item Các công cụ ghi chú hoặc lịch cá nhân để lưu thông tin vận hành phát sinh.
\end{itemize}

\paragraph{3. Trình độ công nghệ thông tin của người sử dụng}
Giáo viên có khả năng sử dụng tốt các công cụ web phổ biến, bảng tính và
Discord. Tuy nhiên, việc phối hợp nhiều công cụ làm tăng khối lượng thao tác và
tạo nhu cầu về một hệ thống quản lý tập trung hơn.

\subsection{Đánh giá hiện trạng và những tính năng mong muốn}

\subsubsection{Ưu điểm}

\begin{itemize}[leftmargin=*]
\item \textbf{Chi phí ban đầu thấp:} Các công cụ như Google Sheets, Discord và Codeforces dễ tiếp cận.
\item \textbf{Dễ bắt đầu sử dụng:} Giáo viên đã quen với các công cụ hiện tại.
\item \textbf{Linh hoạt:} Có thể thay đổi nhanh nội dung ghi chú, lịch học hoặc thông báo.
\end{itemize}

\subsubsection{Nhược điểm và khó khăn cần giải quyết}

\paragraph{a) Dữ liệu phân tán}
Hồ sơ học sinh, lớp học, điểm danh, tài chính, Discord và Codeforces Gym chưa
được liên kết thành một hồ sơ thống nhất.

\paragraph{b) Đối soát tài chính phức tạp}
Khoản phí phát sinh theo buổi học và giao dịch thu/hoàn tiền chưa được quản lý
tập trung, khiến việc tính số dư học sinh dễ sai lệch.

\paragraph{c) Thiếu lịch sử ghi danh và điểm danh rõ ràng}
Giáo viên khó truy vết học sinh đã học lớp nào, nghỉ từ khi nào, đã được điểm
danh ra sao và khoản phí nào phát sinh từ buổi học nào.

\paragraph{d) Thông báo Discord còn thủ công}
Việc gửi tin nhắn vào channel hoặc DM cho học sinh phụ thuộc vào thao tác thủ
công, dễ gửi sai hoặc bỏ sót người nhận.

\paragraph{e) Theo dõi Codeforces Gym chưa tập trung}
Kết quả Gym chưa được hiển thị theo lớp và theo danh sách học sinh đang học, làm
giảm hiệu quả theo dõi tiến độ luyện tập.

\subsubsection{Giải pháp và tính năng mong muốn}

\begin{table}[!htbp]
\centering
\tablesetup
\begin{tabularx}{\textwidth}{|L{3cm}|Y|}
\hline
\textbf{Chủ thể} & \textbf{Mong muốn} \\
\hline
Giáo viên & Có một hệ thống duy nhất để quản lý học sinh, ghi danh, lớp học, lịch học, buổi học, điểm danh, công nợ và báo cáo. \\
\hline
Giáo viên & Tự động liên kết điểm danh, buổi học và khoản phí để tính số dư học sinh chính xác hơn. \\
\hline
Giáo viên & Gắn Codeforces Gym vào lớp học và xem standing theo danh sách học sinh đang học. \\
\hline
Giáo viên & Cấu hình Discord guild/channel cho lớp và gửi thông báo hoặc tin nhắn riêng từ hệ thống. \\
\hline
Quản trị viên & Quản lý tài khoản giáo viên và cấu hình bot Discord dùng chung một cách bảo mật. \\
\hline
Học sinh & Uỷ quyền Discord một lần để được hệ thống hỗ trợ tham gia guild lớp và nhận thông báo phù hợp. \\
\hline
\end{tabularx}
\end{table}

\subsection{Xác định và thu thập yêu cầu}

Từ kết quả khảo sát hiện trạng và các nhu cầu thực tế, nhóm xác định hệ thống
cần tập trung vào các nhóm yêu cầu chính sau:

\begin{itemize}[leftmargin=*]
\item \textbf{Quản lý tài khoản:} Hỗ trợ đăng ký, đăng nhập, cập nhật hồ sơ giáo viên và quản trị tài khoản giáo viên.
\item \textbf{Quản lý học sinh và ghi danh:} Lưu trữ hồ sơ học sinh, chuyển lớp, rút học sinh, khôi phục và lưu trữ học sinh.
\item \textbf{Quản lý lớp học, lịch học và buổi học:} Tạo lớp, thiết lập lịch định kỳ, tạo/hủy buổi học và xem lịch sử buổi học.
\item \textbf{Quản lý điểm danh:} Cập nhật trạng thái điểm danh theo từng học sinh trong buổi học.
\item \textbf{Quản lý tài chính và công nợ:} Quản lý \code{FeeRecord}, \code{Transaction}, số dư học sinh và báo cáo thu nhập.
\item \textbf{Quản lý Discord:} Cấu hình bot, liên kết Discord của giáo viên, uỷ quyền Discord của học sinh, chọn guild/channel và gửi thông báo.
\item \textbf{Quản lý Codeforces Gym:} Gắn Gym cho lớp, đồng bộ \code{GymProblem} và xem bảng xếp hạng Gym.
\item \textbf{Dashboard và báo cáo:} Tổng hợp số liệu lớp học, học sinh, tài chính và hồ sơ học tập dựa trên buổi học, điểm danh, giao dịch và Codeforces Gym.
\end{itemize}

Kế hoạch thu thập yêu cầu được thực hiện thông qua phỏng vấn giáo viên, quan sát
quy trình quản lý hiện tại, nghiên cứu các tài liệu đang được sử dụng và chuẩn
hóa chúng thành các nhóm yêu cầu chức năng, phi chức năng, yêu cầu hệ thống và
yêu cầu công nghệ cho phần mềm mới.

\section{Đặc tả và phân tích yêu cầu chức năng phần mềm}

\subsection{Yêu cầu chức năng}

\subsubsection{\textit{Yêu cầu chức năng nghiệp vụ}}

Hệ thống giúp giáo viên quản lý lớp học lập trình thi đấu trong phiên bản hiện tại là ứng dụng web
quản lý hoạt động dạy thêm thuật toán. Phần mềm hỗ trợ giáo viên quản lý tài
khoản, học sinh, lớp học, lịch học, buổi học, điểm danh, ghi danh, tài chính,
báo cáo, Discord và Codeforces Gym. Ngoài giáo viên, hệ thống có thêm vai trò
quản trị viên để quản lý tài khoản giáo viên và cấu hình bot Discord dùng chung.

Bảng yêu cầu nghiệp vụ:

{\scriptsize
\setlength{\tabcolsep}{2pt}
\renewcommand{\arraystretch}{1.12}
\begin{longtable}{|C{0.7cm}|L{2.5cm}|L{6.4cm}|C{1.05cm}|C{1.05cm}|L{1.8cm}|}
\hline
\textbf{STT} & \textbf{Nghiệp vụ} & \textbf{Thao tác cụ thể} & \textbf{Biểu mẫu} & \textbf{Quy định} & \textbf{Ghi chú} \\
\hline
\endfirsthead
\hline
\textbf{STT} & \textbf{Nghiệp vụ} & \textbf{Thao tác cụ thể} & \textbf{Biểu mẫu} & \textbf{Quy định} & \textbf{Ghi chú} \\
\hline
\endhead
1 & \multirow{4}{2.5cm}{\RaggedRight Quản lý tài khoản} & Đăng nhập hệ thống & BM 1.1 & QĐ 1.1 & Truy cập \\
\cline{1-1}\cline{3-6}
2 & & Đăng ký tài khoản giáo viên & BM 1.2 & QĐ 1.2 & Khởi tạo \\
\cline{1-1}\cline{3-6}
3 & & Cập nhật hồ sơ và thông tin Codeforces của giáo viên & BM 1.3 & QĐ 1.3 & Cá nhân \\
\cline{1-1}\cline{3-6}
4 & & Liên kết tài khoản Discord của giáo viên & BM 1.4 & QĐ 1.4 & OAuth \\
\hline
5 & \multirow{5}{2.5cm}{\RaggedRight Quản lý học sinh và ghi danh} & Thêm học sinh và ghi danh vào lớp & BM 1.5 & QĐ 1.5 & Quản lý \\
\cline{1-1}\cline{3-6}
6 & & Cập nhật hồ sơ học sinh & BM 1.6 & QĐ 1.6 & Quản lý \\
\cline{1-1}\cline{3-6}
7 & & Chuyển lớp, rút học sinh, khôi phục học sinh & BM 1.7 & QĐ 1.7 & Ghi danh \\
\cline{1-1}\cline{3-6}
8 & & Lưu trữ học sinh đang chờ xử lý công nợ & BM 1.8 & QĐ 1.8 & Lưu trữ \\
\cline{1-1}\cline{3-6}
9 & & Tra cứu và xem hồ sơ học tập của học sinh &  &  & Tra cứu \\
\hline
10 & \multirow{4}{2.5cm}{\RaggedRight Quản lý lớp học, lịch học và buổi học} & Thêm/sửa/lưu trữ lớp học & BM 1.10 & QĐ 1.10 & Lớp học \\
\cline{1-1}\cline{3-6}
11 & & Thiết lập lịch học định kỳ của lớp & BM 1.11 & QĐ 1.11 & Lịch học \\
\cline{1-1}\cline{3-6}
12 & & Tạo buổi học thủ công và xem danh sách buổi học & BM 1.12 & QĐ 1.12 & Buổi học \\
\cline{1-1}\cline{3-6}
13 & & Hủy buổi học &  & QĐ 1.13 & Vận hành \\
\hline
14 & Quản lý điểm danh & Xem và cập nhật trạng thái điểm danh học sinh theo buổi học & BM 1.14 & QĐ 1.14 & Điểm danh \\
\hline
15 & \multirow{4}{2.5cm}{\RaggedRight Quản lý tài chính và công nợ} & Ghi nhận giao dịch thu tiền hoặc hoàn tiền & BM 1.15 & QĐ 1.15 & Giao dịch \\
\cline{1-1}\cline{3-6}
16 & & Cập nhật trạng thái khoản phí phát sinh từ buổi học & BM 1.16 & QĐ 1.16 & Khoản phí \\
\cline{1-1}\cline{3-6}
17 & & Xem số dư học sinh và tổng quan tài chính &  &  & Đối soát \\
\cline{1-1}\cline{3-6}
18 & & Lập báo cáo thu nhập theo kỳ & BM 1.18 & QĐ 1.18 & Báo cáo \\
\hline
19 & \multirow{4}{2.5cm}{\RaggedRight Quản lý Discord} & Cấu hình bot Discord toàn hệ thống & BM 1.19 & QĐ 1.19 & Quản trị \\
\cline{1-1}\cline{3-6}
20 & & Chọn Discord guild và channel cho lớp & BM 1.20 & QĐ 1.20 & Liên kết \\
\cline{1-1}\cline{3-6}
21 & & Gửi thông báo vào channel của lớp & BM 1.21 & QĐ 1.21 & Gửi tin \\
\cline{1-1}\cline{3-6}
22 & & Gửi tin nhắn riêng cho học sinh qua Discord & BM 1.22 & QĐ 1.22 & DM \\
\hline
23 & \multirow{2}{2.5cm}{\RaggedRight Quản lý Codeforces Gym} & Gắn hoặc gỡ Gym Codeforces khỏi lớp học & BM 1.23 & QĐ 1.23 & Tích hợp \\
\cline{1-1}\cline{3-6}
24 & & Xem bảng xếp hạng Gym và hồ sơ học tập &  &  & Theo dõi \\
\hline
25 & Uỷ quyền Discord của học sinh & Học sinh uỷ quyền Discord OAuth để hệ thống lưu \code{StudentDiscordCredential} và quản lý membership Discord của lớp & BM 1.24 & QĐ 1.24 & OAuth \\
\hline
26 & \multirow{2}{2.5cm}{\RaggedRight Quản trị hệ thống} & Quản lý tài khoản giáo viên & BM 1.25 & QĐ 1.25 & Admin \\
\cline{1-1}\cline{3-6}
27 & & Xem tổng quan giảng dạy &  &  & Tổng quan \\
\hline
\end{longtable}
}

\paragraph{Biểu mẫu 1 và quy định 1}

\textbf{* BM 1.1: Đăng nhập hệ thống}

\begin{table}[!htbp]
\centering
\tablesetup
\begin{tabularx}{0.78\textwidth}{|Y|Y|}
\hline
\textbf{Tên đăng nhập} & \textbf{Mật khẩu} \\
\hline
\end{tabularx}
\end{table}

\noindent - QĐ 1.1: Tên đăng nhập và mật khẩu không được để trống; tài khoản phải tồn tại, mật khẩu phải khớp, tài khoản phải đang hoạt động. Sau khi đăng nhập thành công, hệ thống cấp token và phân quyền theo vai trò giáo viên hoặc quản trị viên.

\paragraph{Biểu mẫu 2 và quy định 2}

\textbf{* BM 1.2: Đăng ký tài khoản giáo viên}

\begin{table}[!htbp]
\centering
\tablesetup
\begin{tabularx}{\textwidth}{|C{1cm}|L{3.2cm}|L{3cm}|Y|}
\hline
\textbf{STT} & \textbf{Trường/Mục} & \textbf{Cách nhập/hiển thị} & \textbf{Ghi chú} \\
\hline
1 & Tên đăng nhập & Nhập văn bản & Bắt buộc, dùng để đăng nhập hệ thống \\
\hline
2 & Mật khẩu & Nhập bảo mật & Bắt buộc \\
\hline
3 & Codeforces handle & Nhập văn bản & Không bắt buộc, dùng cho đồng bộ Gym \\
\hline
4 & Codeforces API key & Nhập bảo mật & Không bắt buộc \\
\hline
5 & Codeforces API secret & Nhập bảo mật & Không bắt buộc \\
\hline
\end{tabularx}
\end{table}

\noindent - QĐ 1.2: Tên đăng nhập và mật khẩu là bắt buộc; tên đăng nhập không được trùng; các thông tin Codeforces có thể bỏ trống nhưng nếu nhập thì phải được lưu bảo mật và chỉ sử dụng cho việc đồng bộ dữ liệu Gym.

\paragraph{Biểu mẫu 3 và quy định 3}

\textbf{* BM 1.3: Cập nhật hồ sơ giáo viên}

\begin{table}[!htbp]
\centering
\tablesetup
\begin{tabularx}{\textwidth}{|C{1cm}|L{3.2cm}|L{3cm}|Y|}
\hline
\textbf{STT} & \textbf{Trường/Mục} & \textbf{Cách nhập/hiển thị} & \textbf{Ghi chú} \\
\hline
1 & Tên đăng nhập & Nhập văn bản & Có thể cập nhật \\
\hline
2 & Mật khẩu mới & Nhập bảo mật & Không bắt buộc \\
\hline
3 & Codeforces handle & Nhập văn bản & Có thể xóa hoặc cập nhật \\
\hline
4 & Codeforces API key & Nhập bảo mật & Không hiển thị đầy đủ sau khi lưu \\
\hline
5 & Codeforces API secret & Nhập bảo mật & Không hiển thị đầy đủ sau khi lưu \\
\hline
\end{tabularx}
\end{table}

\noindent - QĐ 1.3: Khi cập nhật hồ sơ phải có ít nhất một trường thay đổi; mật khẩu không được rỗng nếu được gửi lên; thông tin nhạy cảm Codeforces phải được bảo vệ khi lưu trữ và không được hiển thị đầy đủ ở giao diện danh sách.

\paragraph{Biểu mẫu 4 và quy định 4}

\textbf{* BM 1.4: Liên kết Discord của giáo viên}

\begin{table}[!htbp]
\centering
\tablesetup
\begin{tabularx}{\textwidth}{|C{1cm}|L{3.2cm}|L{3cm}|Y|}
\hline
\textbf{STT} & \textbf{Trường/Mục} & \textbf{Cách nhập/hiển thị} & \textbf{Ghi chú} \\
\hline
1 & Yêu cầu xác minh Discord & Nút thao tác & Hệ thống sinh URL OAuth \\
\hline
2 & Discord user id & Tự động ghi nhận & Ghi sau khi OAuth thành công \\
\hline
3 & Discord username & Tự động hiển thị & Dùng để nhận diện giáo viên \\
\hline
\end{tabularx}
\end{table}

\noindent - QĐ 1.4: Giáo viên phải đăng nhập trước khi bắt đầu xác minh; callback OAuth phải hợp lệ; hệ thống chỉ ghi nhận Discord identity khi Discord trả về thông tin xác thực hợp lệ.

\paragraph{Biểu mẫu 5 và quy định 5}

\textbf{* BM 1.5: Thêm học sinh và ghi danh vào lớp}

\begin{table}[!htbp]
\centering
\tablesetup
\begin{tabularx}{\textwidth}{|C{1cm}|L{3.2cm}|L{3cm}|Y|}
\hline
\textbf{STT} & \textbf{Trường/Mục} & \textbf{Cách nhập/hiển thị} & \textbf{Ghi chú} \\
\hline
1 & Họ và tên học sinh & Nhập văn bản & Bắt buộc \\
\hline
2 & Lớp học & Chọn từ danh sách & Bắt buộc, lớp thuộc giáo viên hiện tại \\
\hline
3 & Codeforces handle & Nhập văn bản & Bắt buộc \\
\hline
4 & Số điện thoại & Nhập văn bản & Không bắt buộc \\
\hline
5 & Ghi chú & Nhập nội dung dài & Không bắt buộc \\
\hline
6 & Ngày ghi danh & Chọn ngày giờ & Mặc định là thời điểm hiện tại \\
\hline
\end{tabularx}
\end{table}

\noindent - QĐ 1.5: Họ tên, lớp học và Codeforces handle không được để trống; lớp học phải thuộc giáo viên đang thao tác; mỗi học sinh chỉ có tối đa một ghi danh đang hoạt động trong phạm vi một giáo viên.

\paragraph{Biểu mẫu 6 và quy định 6}

\textbf{* BM 1.6: Cập nhật hồ sơ học sinh}

\begin{table}[!htbp]
\centering
\tablesetup
\begin{tabularx}{\textwidth}{|C{1cm}|L{3.2cm}|L{3cm}|Y|}
\hline
\textbf{STT} & \textbf{Trường/Mục} & \textbf{Cách nhập/hiển thị} & \textbf{Ghi chú} \\
\hline
1 & Họ và tên học sinh & Nhập văn bản & Có thể cập nhật \\
\hline
2 & Codeforces handle & Nhập văn bản & Có thể cập nhật \\
\hline
3 & Số điện thoại & Nhập văn bản & Có thể để trống \\
\hline
4 & Ghi chú & Nhập nội dung dài & Có thể để trống \\
\hline
5 & Discord username & Tự động hiển thị & Sinh từ liên kết Discord của học sinh \\
\hline
6 & Trạng thái học sinh & Tự động hiển thị & Đang học, chờ lưu trữ hoặc đã lưu trữ \\
\hline
\end{tabularx}
\end{table}

\noindent - QĐ 1.6: Khi cập nhật phải có ít nhất một trường thay đổi; học sinh phải thuộc giáo viên hiện tại; nếu cập nhật Codeforces handle thì không được để rỗng.

\paragraph{Biểu mẫu 7 và quy định 7}

\textbf{* BM 1.7: Chuyển lớp, rút học sinh và khôi phục học sinh}

\begin{table}[!htbp]
\centering
\tablesetup
\begin{tabularx}{\textwidth}{|C{1cm}|L{3.2cm}|L{3cm}|Y|}
\hline
\textbf{STT} & \textbf{Trường/Mục} & \textbf{Cách nhập/hiển thị} & \textbf{Ghi chú} \\
\hline
1 & Học sinh & Chọn từ danh sách & Bắt buộc \\
\hline
2 & Lớp chuyển đến & Chọn từ danh sách & Dùng khi chuyển lớp hoặc khôi phục \\
\hline
3 & Thời điểm chuyển/rút/khôi phục & Chọn ngày giờ & Mặc định là hiện tại \\
\hline
\end{tabularx}
\end{table}

\noindent - QĐ 1.7: Học sinh và lớp chuyển đến phải thuộc giáo viên hiện tại; khi chuyển lớp, hệ thống kết thúc ghi danh cũ và tạo ghi danh mới; khi rút học sinh, hệ thống xác định trạng thái tiếp theo dựa trên số dư tài chính còn phải thu hoặc cần hoàn.

\paragraph{Biểu mẫu 8 và quy định 8}

\textbf{* BM 1.8: Lưu trữ học sinh chờ xử lý}

\begin{table}[!htbp]
\centering
\tablesetup
\begin{tabularx}{\textwidth}{|C{1cm}|L{3.2cm}|L{3cm}|Y|}
\hline
\textbf{STT} & \textbf{Trường/Mục} & \textbf{Cách nhập/hiển thị} & \textbf{Ghi chú} \\
\hline
1 & Học sinh & Chọn từ danh sách & Bắt buộc \\
\hline
2 & Thời điểm lưu trữ & Chọn ngày giờ & Mặc định là hiện tại \\
\hline
3 & Lý do chờ lưu trữ & Tự động hiển thị & Cần thu thêm hoặc cần hoàn tiền \\
\hline
\end{tabularx}
\end{table}

\noindent - QĐ 1.8: Chỉ cho phép lưu trữ khi học sinh đang ở trạng thái chờ lưu trữ; nếu còn số dư cần xử lý thì hệ thống phải thể hiện lý do để giáo viên đối soát trước.

\paragraph{Biểu mẫu 10 và quy định 10}

\textbf{* BM 1.10: Thêm/sửa/lưu trữ lớp học}

\begin{table}[!htbp]
\centering
\tablesetup
\begin{tabularx}{\textwidth}{|C{1cm}|L{3.2cm}|L{3cm}|Y|}
\hline
\textbf{STT} & \textbf{Trường/Mục} & \textbf{Cách nhập/hiển thị} & \textbf{Ghi chú} \\
\hline
1 & Tên lớp & Nhập văn bản & Bắt buộc \\
\hline
2 & Học phí mỗi buổi & Nhập số & Bắt buộc, không âm \\
\hline
3 & Danh sách lịch học & Nhập danh sách & Ít nhất một lịch học \\
\hline
4 & Trạng thái lớp & Tự động hoặc chọn thao tác & Đang hoạt động hoặc đã lưu trữ \\
\hline
\end{tabularx}
\end{table}

\noindent - QĐ 1.10: Tên lớp không được để trống; học phí mỗi buổi là số nguyên không âm; lớp phải có ít nhất một lịch học khi tạo mới; khi lưu trữ lớp, hệ thống chuyển trạng thái lớp sang đã lưu trữ thay vì xóa dữ liệu.

\paragraph{Biểu mẫu 11 và quy định 11}

\textbf{* BM 1.11: Thiết lập lịch học định kỳ}

\begin{table}[!htbp]
\centering
\tablesetup
\begin{tabularx}{\textwidth}{|C{1cm}|L{3.2cm}|L{3cm}|Y|}
\hline
\textbf{STT} & \textbf{Trường/Mục} & \textbf{Cách nhập/hiển thị} & \textbf{Ghi chú} \\
\hline
1 & Thứ trong tuần & Chọn từ danh sách & Giá trị 0 đến 6 tương ứng các ngày trong tuần \\
\hline
2 & Giờ bắt đầu & Chọn giờ & Định dạng HH:mm hoặc HH:mm:ss \\
\hline
3 & Giờ kết thúc & Chọn giờ & Phải sau giờ bắt đầu \\
\hline
\end{tabularx}
\end{table}

\noindent - QĐ 1.11: Thứ trong tuần nằm trong miền giá trị hợp lệ; giờ kết thúc phải lớn hơn giờ bắt đầu; khi cập nhật lịch của lớp, danh sách lịch mới thay thế cấu hình lịch đang áp dụng.

\paragraph{Biểu mẫu 12 và quy định 12}

\textbf{* BM 1.12: Tạo buổi học thủ công}

\begin{table}[!htbp]
\centering
\tablesetup
\begin{tabularx}{\textwidth}{|C{1cm}|L{3.2cm}|L{3cm}|Y|}
\hline
\textbf{STT} & \textbf{Trường/Mục} & \textbf{Cách nhập/hiển thị} & \textbf{Ghi chú} \\
\hline
1 & Lớp học & Chọn từ danh sách & Bắt buộc \\
\hline
2 & Ngày học & Chọn ngày & Dùng khi tạo theo ngày và giờ \\
\hline
3 & Giờ bắt đầu & Chọn giờ & Dùng khi tạo theo ngày và giờ \\
\hline
4 & Thời điểm bắt đầu & Chọn ngày giờ & Có thể dùng thay cho ngày và giờ bắt đầu \\
\hline
5 & Giờ kết thúc & Chọn giờ & Bắt buộc \\
\hline
\end{tabularx}
\end{table}

\noindent - QĐ 1.12: Buổi học phải thuộc lớp của giáo viên hiện tại; giáo viên phải nhập một trong hai nhóm dữ liệu: thời điểm bắt đầu và giờ kết thúc, hoặc ngày học, giờ bắt đầu và giờ kết thúc; giờ kết thúc phải sau giờ bắt đầu.

\paragraph{Quy định 13}

\noindent - QĐ 1.13: Chỉ được hủy buổi học thuộc giáo viên hiện tại; buổi học bị hủy phải được giữ lại lịch sử với trạng thái hủy, không xóa cứng khỏi hệ thống; dữ liệu điểm danh và khoản phí liên quan phải được xử lý theo trạng thái buổi học.

\paragraph{Biểu mẫu 14 và quy định 14}

\textbf{* BM 1.14: Cập nhật điểm danh buổi học}

\begin{table}[!htbp]
\centering
\tablesetup
\begin{tabularx}{\textwidth}{|C{1cm}|L{3.2cm}|L{3cm}|Y|}
\hline
\textbf{STT} & \textbf{Trường/Mục} & \textbf{Cách nhập/hiển thị} & \textbf{Ghi chú} \\
\hline
1 & Buổi học & Chọn từ danh sách & Bắt buộc \\
\hline
2 & Học sinh & Chọn từ danh sách & Bắt buộc \\
\hline
3 & Trạng thái điểm danh & Chọn từ danh sách & Có mặt, vắng có phép hoặc vắng không phép \\
\hline
4 & Ghi chú & Nhập văn bản & Không bắt buộc \\
\hline
5 & Nguồn điểm danh & Tự động hiển thị & Bot, thủ công hoặc hệ thống \\
\hline
\end{tabularx}
\end{table}

\noindent - QĐ 1.14: Trạng thái điểm danh chỉ nhận một trong ba giá trị hợp lệ: có mặt, vắng có phép, vắng không phép; học sinh và buổi học phải thuộc giáo viên hiện tại; nếu điểm danh làm phát sinh hoặc thay đổi khoản phí, hệ thống phải bảo toàn liên kết giữa \code{Attendance} và \code{FeeRecord}.

\paragraph{Biểu mẫu 15 và quy định 15}

\textbf{* BM 1.15: Ghi nhận giao dịch tài chính}

\begin{table}[!htbp]
\centering
\tablesetup
\begin{tabularx}{\textwidth}{|C{1cm}|L{3.2cm}|L{3cm}|Y|}
\hline
\textbf{STT} & \textbf{Trường/Mục} & \textbf{Cách nhập/hiển thị} & \textbf{Ghi chú} \\
\hline
1 & Học sinh & Chọn từ danh sách & Bắt buộc \\
\hline
2 & Số tiền & Nhập số & Bắt buộc \\
\hline
3 & Loại giao dịch & Chọn từ danh sách & Thu tiền hoặc hoàn tiền \\
\hline
4 & Ghi chú & Nhập nội dung & Không bắt buộc \\
\hline
5 & Thời điểm ghi nhận & Chọn ngày giờ & Mặc định là hiện tại \\
\hline
\end{tabularx}
\end{table}

\noindent - QĐ 1.15: Học sinh phải thuộc giáo viên hiện tại; loại giao dịch chỉ được là thu tiền hoặc hoàn tiền; giao dịch thu tiền lưu số tiền dương, giao dịch hoàn tiền lưu số tiền âm; khi sửa giao dịch đã có, hệ thống phải kiểm tra lại quyền sở hữu học sinh và giao dịch.

\paragraph{Biểu mẫu 16 và quy định 16}

\textbf{* BM 1.16: Cập nhật trạng thái khoản phí}

\begin{table}[!htbp]
\centering
\tablesetup
\begin{tabularx}{\textwidth}{|C{1cm}|L{3.2cm}|L{3cm}|Y|}
\hline
\textbf{STT} & \textbf{Trường/Mục} & \textbf{Cách nhập/hiển thị} & \textbf{Ghi chú} \\
\hline
1 & Khoản phí & Chọn từ danh sách & Bắt buộc \\
\hline
2 & Trạng thái khoản phí & Chọn từ danh sách & Đang hiệu lực hoặc đã hủy \\
\hline
\end{tabularx}
\end{table}

\noindent - QĐ 1.16: Khoản phí phải thuộc học sinh hoặc buổi học của giáo viên hiện tại; trạng thái khoản phí chỉ được là đang hiệu lực hoặc đã hủy; thao tác hủy khoản phí không làm mất lịch sử phát sinh.

\paragraph{Biểu mẫu 18 và quy định 18}

\textbf{* BM 1.18: Lập báo cáo thu nhập}

\begin{table}[!htbp]
\centering
\tablesetup
\begin{tabularx}{\textwidth}{|C{1cm}|L{3.2cm}|L{3cm}|Y|}
\hline
\textbf{STT} & \textbf{Trường/Mục} & \textbf{Cách nhập/hiển thị} & \textbf{Ghi chú} \\
\hline
1 & Từ ngày & Chọn ngày giờ & Không bắt buộc \\
\hline
2 & Đến ngày & Chọn ngày giờ & Không bắt buộc \\
\hline
3 & Danh sách lớp & Chọn nhiều lớp & Không bắt buộc \\
\hline
4 & Bao gồm khoản chưa thu & Bật/tắt & Không bắt buộc \\
\hline
\end{tabularx}
\end{table}

\noindent - QĐ 1.18: Nếu có cả từ ngày và đến ngày thì từ ngày không được sau đến ngày; danh sách lớp được lọc phải thuộc giáo viên hiện tại; báo cáo phải tổng hợp giao dịch, khoản phí và số dư theo đúng kỳ lọc.

\paragraph{Biểu mẫu 19 và quy định 19}

\textbf{* BM 1.19: Cấu hình bot Discord}

\begin{table}[!htbp]
\centering
\tablesetup
\begin{tabularx}{\textwidth}{|C{1cm}|L{3.2cm}|L{3cm}|Y|}
\hline
\textbf{STT} & \textbf{Trường/Mục} & \textbf{Cách nhập/hiển thị} & \textbf{Ghi chú} \\
\hline
1 & Bot token & Nhập bảo mật & Bắt buộc \\
\hline
2 & Client ID & Nhập văn bản & Bắt buộc \\
\hline
3 & Client secret & Nhập bảo mật & Bắt buộc \\
\hline
4 & Permissions & Nhập văn bản & Không bắt buộc \\
\hline
5 & Scopes & Nhập văn bản & Không bắt buộc \\
\hline
\end{tabularx}
\end{table}

\noindent - QĐ 1.19: Chỉ quản trị viên được cấu hình bot Discord; bot token và client secret là dữ liệu nhạy cảm, phải được bảo vệ và không hiển thị đầy đủ sau khi lưu.

\paragraph{Biểu mẫu 20 và quy định 20}

\textbf{* BM 1.20: Chọn Discord guild và channel cho lớp}

\begin{table}[!htbp]
\centering
\tablesetup
\begin{tabularx}{\textwidth}{|C{1cm}|L{3.2cm}|L{3cm}|Y|}
\hline
\textbf{STT} & \textbf{Trường/Mục} & \textbf{Cách nhập/hiển thị} & \textbf{Ghi chú} \\
\hline
1 & Lớp học & Chọn từ danh sách & Bắt buộc \\
\hline
2 & Discord guild & Chọn từ danh sách & Bắt buộc \\
\hline
3 & Channel thông báo & Chọn từ danh sách & Không bắt buộc \\
\hline
4 & Voice channel điểm danh & Chọn từ danh sách & Không bắt buộc \\
\hline
\end{tabularx}
\end{table}

\noindent - QĐ 1.20: Lớp học phải thuộc giáo viên hiện tại; guild và channel phải nằm trong dữ liệu Discord mà hệ thống nhìn thấy; channel thông báo và voice channel điểm danh phải thuộc đúng guild đã chọn.

\paragraph{Biểu mẫu 21 và quy định 21}

\textbf{* BM 1.21: Gửi thông báo channel Discord}

\begin{table}[!htbp]
\centering
\tablesetup
\begin{tabularx}{\textwidth}{|C{1cm}|L{3.2cm}|L{3cm}|Y|}
\hline
\textbf{STT} & \textbf{Trường/Mục} & \textbf{Cách nhập/hiển thị} & \textbf{Ghi chú} \\
\hline
1 & Nội dung thông báo & Nhập văn bản & Bắt buộc \\
\hline
2 & Danh sách guild nhận thông báo & Chọn nhiều guild & Ít nhất một guild \\
\hline
\end{tabularx}
\end{table}

\noindent - QĐ 1.21: Nội dung không được để trống; danh sách guild phải có ít nhất một phần tử; hệ thống loại bỏ guild trùng trong một lần gửi; khi gửi lỗi, hệ thống cần trả trạng thái để giáo viên biết thao tác không thành công.

\paragraph{Biểu mẫu 22 và quy định 22}

\textbf{* BM 1.22: Gửi tin nhắn riêng cho học sinh}

\begin{table}[!htbp]
\centering
\tablesetup
\begin{tabularx}{\textwidth}{|C{1cm}|L{3.2cm}|L{3cm}|Y|}
\hline
\textbf{STT} & \textbf{Trường/Mục} & \textbf{Cách nhập/hiển thị} & \textbf{Ghi chú} \\
\hline
1 & Nội dung tin nhắn & Nhập văn bản & Bắt buộc \\
\hline
2 & Học sinh nhận tin & Chọn một hoặc nhiều học sinh & Dùng khi gửi cho danh sách học sinh \\
\hline
3 & Lớp nhận tin & Chọn lớp & Dùng khi gửi cho toàn bộ học sinh trong lớp \\
\hline
\end{tabularx}
\end{table}

\noindent - QĐ 1.22: Nội dung không được để trống; phải chọn ít nhất một học sinh hoặc một lớp; học sinh và lớp phải thuộc giáo viên hiện tại; hệ thống chỉ gửi được cho học sinh đã có thông tin Discord phù hợp.

\paragraph{Biểu mẫu 23 và quy định 23}

\textbf{* BM 1.23: Gắn Codeforces Gym cho lớp}

\begin{table}[!htbp]
\centering
\tablesetup
\begin{tabularx}{\textwidth}{|C{1cm}|L{3.2cm}|L{3cm}|Y|}
\hline
\textbf{STT} & \textbf{Trường/Mục} & \textbf{Cách nhập/hiển thị} & \textbf{Ghi chú} \\
\hline
1 & Lớp học & Chọn từ danh sách & Bắt buộc \\
\hline
2 & Gym ID & Nhập số hoặc chọn từ danh sách & Bắt buộc, phải là chuỗi số \\
\hline
3 & Chu kỳ đồng bộ & Nhập số phút & Không bắt buộc, số nguyên dương \\
\hline
\end{tabularx}
\end{table}

\noindent - QĐ 1.23: Gym ID phải là số; lớp học phải thuộc giáo viên hiện tại; chu kỳ đồng bộ nếu nhập thì phải là số nguyên dương; khi gỡ Gym khỏi lớp, hệ thống phải giữ dữ liệu lịch sử đã đồng bộ nếu cần phục vụ báo cáo học tập.

\paragraph{Biểu mẫu 24 và quy định 24}

\textbf{* BM 1.24: Uỷ quyền Discord của học sinh}

\begin{table}[!htbp]
\centering
\tablesetup
\begin{tabularx}{\textwidth}{|C{1cm}|L{3.2cm}|L{3cm}|Y|}
\hline
\textbf{STT} & \textbf{Trường/Mục} & \textbf{Cách nhập/hiển thị} & \textbf{Ghi chú} \\
\hline
1 & Học sinh & Xác định từ liên kết uỷ quyền & Bắt buộc \\
\hline
2 & Discord OAuth consent & Học sinh xác nhận trên Discord & Dùng quyền nhận diện và tham gia guild \\
\hline
3 & Discord user id & Tự động ghi nhận & Lưu vào \code{StudentDiscordCredential} \\
\hline
4 & Access token & Tự động ghi nhận bảo mật & Dùng để quản lý membership theo lớp \\
\hline
\end{tabularx}
\end{table}

\noindent - QĐ 1.24: URL uỷ quyền phải gắn với học sinh hợp lệ; callback Discord phải hợp lệ; hệ thống chỉ lưu Discord credential khi OAuth thành công; token của học sinh là dữ liệu nhạy cảm và không được hiển thị đầy đủ.

\paragraph{Biểu mẫu 25 và quy định 25}

\textbf{* BM 1.25: Quản lý tài khoản giáo viên}

\begin{table}[!htbp]
\centering
\tablesetup
\begin{tabularx}{\textwidth}{|C{1cm}|L{3.2cm}|L{3cm}|Y|}
\hline
\textbf{STT} & \textbf{Trường/Mục} & \textbf{Cách nhập/hiển thị} & \textbf{Ghi chú} \\
\hline
1 & Giáo viên & Chọn từ danh sách & Bắt buộc khi cập nhật \\
\hline
2 & Tên đăng nhập & Nhập văn bản & Có thể cập nhật \\
\hline
3 & Mật khẩu & Nhập bảo mật & Có thể đặt lại \\
\hline
4 & Trạng thái hoạt động & Bật/tắt & Kích hoạt hoặc vô hiệu hóa \\
\hline
5 & Codeforces credential & Nhập hoặc cập nhật & Không bắt buộc \\
\hline
\end{tabularx}
\end{table}

\noindent - QĐ 1.25: Chỉ quản trị viên được quản lý tài khoản giáo viên; khi cập nhật phải có ít nhất một trường thay đổi; không xóa cứng tài khoản giáo viên mà chỉ chuyển trạng thái hoạt động khi cần ngừng truy cập.

\clearpage
\subsubsection{Bảng trách nhiệm yêu cầu nghiệp vụ}

\begin{table}[!htbp]
\centering
\scriptsize
\setlength{\tabcolsep}{2pt}
\renewcommand{\arraystretch}{1.2}
\begin{tabularx}{\textwidth}{|C{0.8cm}|L{2.6cm}|L{3.5cm}|Y|L{1.6cm}|}
\hline
\textbf{STT} & \textbf{Tên yêu cầu} & \textbf{Người dùng} & \textbf{Phần mềm} & \textbf{Ghi chú} \\
\hline
1 & Quản lý tài khoản & Giáo viên hoặc quản trị viên nhập thông tin đăng nhập, đăng ký hoặc cập nhật hồ sơ. & Xác thực, cấp token, kiểm tra vai trò, lưu credential Codeforces và hỗ trợ xác minh Discord của giáo viên. & Truy cập hệ thống. \\
\hline
2 & Quản lý học sinh và ghi danh & Giáo viên thêm, sửa, tra cứu, chuyển lớp, rút, khôi phục hoặc lưu trữ học sinh. & Kiểm tra quyền sở hữu dữ liệu, quản lý hồ sơ học sinh, trạng thái ghi danh và trạng thái lưu trữ dựa trên công nợ. & Trung tâm dữ liệu học sinh. \\
\hline
3 & Quản lý lớp học, lịch học và buổi học & Giáo viên tạo lớp, cập nhật lịch định kỳ, tạo buổi học thủ công, xem hoặc hủy buổi học. & Lưu lớp, lịch, buổi học, trạng thái lớp và trạng thái buổi học; đảm bảo dữ liệu thuộc đúng giáo viên. & Vận hành lớp học. \\
\hline
4 & Quản lý điểm danh & Giáo viên cập nhật trạng thái điểm danh từng học sinh trong buổi học. & Lưu trạng thái điểm danh, nguồn điểm danh, ghi chú và liên kết dữ liệu với khoản phí phát sinh. & Phục vụ tính phí. \\
\hline
5 & Quản lý tài chính và công nợ & Giáo viên ghi nhận giao dịch, cập nhật khoản phí, xem số dư và báo cáo. & Quản lý \code{Transaction}, \code{FeeRecord}, số dư học sinh, tổng quan tài chính và báo cáo thu nhập theo kỳ. & Đối soát tài chính. \\
\hline
6 & Quản lý Discord & Quản trị viên cấu hình bot; giáo viên liên kết guild, channel, gửi thông báo và DM. & Lưu cấu hình bot, xử lý OAuth, cache guild/channel, gửi tin nhắn và ghi nhận kết quả. & Phụ thuộc Discord. \\
\hline
7 & Quản lý Codeforces Gym & Giáo viên gắn Gym cho lớp và xem bảng xếp hạng. & Đồng bộ Gym, \code{GymProblem}, \code{GymStanding} và hiển thị ma trận kết quả theo học sinh trong lớp. & Phụ thuộc Codeforces. \\
\hline
8 & Uỷ quyền Discord của học sinh & Học sinh mở liên kết OAuth và xác nhận quyền trên Discord. & Lưu \code{StudentDiscordCredential}, phục vụ mời học sinh vào guild lớp và gửi tin nhắn phù hợp. & Phụ thuộc Discord. \\
\hline
9 & Quản trị hệ thống & Quản trị viên quản lý tài khoản giáo viên và cấu hình bot Discord. & Phân quyền admin, kiểm soát trạng thái tài khoản giáo viên và credential dùng chung. & Vận hành nội bộ. \\
\hline
10 & Dashboard và báo cáo học tập & Giáo viên xem tổng quan và hồ sơ học tập của học sinh. & Tổng hợp số lớp, số học sinh, tài chính, lịch sử Gym và trạng thái học tập. & Theo dõi nhanh. \\
\hline
\end{tabularx}
\end{table}

\subsection{Yêu cầu phi chức năng}

\subsubsection{Yêu cầu hiệu quả}

Yêu cầu hiệu quả liên quan đến tốc độ xử lý và khả năng sử dụng tài nguyên của
hệ thống, đặc biệt là các tác vụ thường xuyên và quan trọng như đăng nhập, tra
cứu học sinh, cập nhật lớp học, điểm danh, ghi nhận tài chính, gửi thông báo và
xem báo cáo.

\noindent - Bảng yêu cầu hiệu quả:

\begin{table}[!htbp]
\centering
\tablesetup
\begin{tabularx}{\textwidth}{|C{0.8cm}|L{2.8cm}|Y|L{2.1cm}|L{1.7cm}|}
\hline
\textbf{STT} & \textbf{Nghiệp vụ} & \textbf{Tốc độ xử lý} & \textbf{Dung lượng lưu trữ} & \textbf{Ghi chú} \\
\hline
1 & Quản lý tài khoản & Đăng nhập và lấy thông tin cá nhân dưới 2 giây; cập nhật hồ sơ dưới 3 giây. & Tăng theo số tài khoản và credential. & Phụ thuộc mã hóa mật khẩu. \\
\hline
2 & Quản lý học sinh và ghi danh & Lưu hồ sơ, chuyển lớp hoặc rút học sinh dưới 5 giây; lọc danh sách dưới 8 giây. & Tăng theo số học sinh và lịch sử ghi danh. & Cần truy vấn số dư. \\
\hline
3 & Quản lý lớp, lịch và buổi học & Tạo/cập nhật lớp dưới 5 giây; danh sách buổi học dưới 8 giây. & Tăng theo lịch học và buổi học phát sinh. & Dữ liệu tăng theo thời gian. \\
\hline
4 & Quản lý điểm danh & Cập nhật một bản ghi điểm danh dưới 3 giây; tải điểm danh của một buổi học dưới 5 giây. & Tăng theo số buổi học và số học sinh. & Có thể liên quan \code{FeeRecord}. \\
\hline
5 & Tài chính và công nợ & Lưu giao dịch dưới 5 giây; tính số dư hoặc tổng quan tài chính dưới 10 giây. & Tăng theo \code{Transaction} và \code{FeeRecord}. & Ưu tiên độ chính xác. \\
\hline
6 & Discord & Lấy guild/channel dưới 10 giây; gửi thông báo hoặc DM dưới 10 giây mỗi lô nhỏ. & Tăng theo credential và cache guild/channel. & Phụ thuộc Discord API. \\
\hline
7 & Codeforces Gym & Xem bảng xếp hạng đã đồng bộ dưới 5 giây; đồng bộ từ Codeforces theo cơ chế nền. & Tăng theo Gym, Problem và Standing. & Phụ thuộc Codeforces. \\
\hline
8 & Báo cáo & Lập báo cáo thu nhập theo kỳ dưới 10 giây với dữ liệu thông thường. & Không phát sinh nhiều nếu chỉ xem trực tiếp. & Cần lọc theo lớp và thời gian. \\
\hline
\end{tabularx}
\end{table}

\noindent - Bảng trách nhiệm yêu cầu hiệu quả:

\begin{table}[!htbp]
\centering
\scriptsize
\setlength{\tabcolsep}{2pt}
\renewcommand{\arraystretch}{1.2}
\begin{tabularx}{\textwidth}{|C{0.8cm}|L{2.7cm}|L{3.4cm}|Y|L{1.5cm}|}
\hline
\textbf{STT} & \textbf{Nghiệp vụ} & \textbf{Người dùng} & \textbf{Phần mềm} & \textbf{Ghi chú} \\
\hline
1 & Tài khoản & Nhập thông tin đăng nhập hoặc cập nhật hồ sơ. & Xác thực, băm mật khẩu, cấp token và trả kết quả trong ngưỡng thời gian. & \\
\hline
2 & Học sinh và ghi danh & Nhập hồ sơ, tiêu chí lọc hoặc yêu cầu thay đổi trạng thái ghi danh. & Tối ưu truy vấn học sinh, ghi danh và số dư liên quan. & \\
\hline
3 & Lớp và buổi học & Tạo lớp, xem lịch, tạo/hủy buổi học. & Lưu lịch học, trạng thái buổi học và truy vấn theo khoảng thời gian hiệu quả. & \\
\hline
4 & Điểm danh & Cập nhật trạng thái học sinh trong buổi học. & Ghi điểm danh nhanh, cập nhật dữ liệu tài chính liên quan khi cần. & \\
\hline
5 & Tài chính & Nhập giao dịch hoặc xem báo cáo. & Tính toán \code{Transaction}, \code{FeeRecord} và số dư nhất quán. & \\
\hline
6 & Tích hợp ngoài & Kích hoạt thao tác Discord hoặc Codeforces. & Dùng dữ liệu cache khi có thể, ghi nhận lỗi khi dịch vụ ngoài phản hồi chậm hoặc thất bại. & \\
\hline
\end{tabularx}
\end{table}

\subsubsection{Yêu cầu tiện dụng}

Yêu cầu tiện dụng tập trung vào khả năng học nhanh, thao tác thuận tiện và hạn
chế sai sót cho giáo viên khi sử dụng hệ thống trong quá trình vận hành lớp học.

\noindent - Bảng yêu cầu tiện dụng:

\begin{table}[!htbp]
\centering
\tablesetup
\begin{tabularx}{\textwidth}{|C{0.9cm}|L{3.1cm}|C{2.3cm}|Y|L{1.7cm}|}
\hline
\textbf{STT} & \textbf{Nghiệp vụ} & \textbf{Mức độ dễ học} & \textbf{Mức độ dễ sử dụng} & \textbf{Ghi chú} \\
\hline
1 & Quản lý tài khoản & 3--5 phút & Màn hình đăng nhập, hồ sơ cá nhân và thông báo lỗi rõ ràng. & \\
\hline
2 & Quản lý học sinh & 10 phút & Dễ tìm kiếm, lọc theo lớp/trạng thái, xem chi tiết và thao tác chuyển/rút/khôi phục. & \\
\hline
3 & Quản lý lớp và buổi học & 15 phút & Form lớp, lịch và buổi học thống nhất, dễ nhận biết trạng thái. & \\
\hline
4 & Điểm danh & 5--10 phút & Cập nhật nhanh theo danh sách học sinh trong một buổi học. & \\
\hline
5 & Tài chính & 10 phút & Số dư, khoản phí và giao dịch hiển thị rõ, hạn chế nhầm lẫn thu/hoàn tiền. & \\
\hline
6 & Discord & 15 phút & Quy trình cấu hình bot, chọn guild/channel và gửi tin cần tuần tự, dễ kiểm tra. & \\
\hline
7 & Codeforces Gym & 5--10 phút & Dễ gắn Gym, xem standing và nhận biết trạng thái dữ liệu. & \\
\hline
8 & Báo cáo & 10 phút & Dễ chọn kỳ, lớp và tùy chọn bao gồm khoản chưa thu. & \\
\hline
\end{tabularx}
\end{table}

\noindent - Bảng trách nhiệm yêu cầu tiện dụng:

\begin{table}[!htbp]
\centering
\scriptsize
\setlength{\tabcolsep}{2pt}
\renewcommand{\arraystretch}{1.2}
\begin{tabularx}{\textwidth}{|C{0.8cm}|L{2.8cm}|L{3.2cm}|Y|L{1.6cm}|}
\hline
\textbf{STT} & \textbf{Nghiệp vụ} & \textbf{Người dùng} & \textbf{Phần mềm} & \textbf{Ghi chú} \\
\hline
1 & Tài khoản & Nhập thông tin và đọc thông báo lỗi. & Kiểm tra dữ liệu tức thời, thông báo lỗi ngắn gọn. & \\
\hline
2 & Học sinh & Tìm kiếm, xem chi tiết và thao tác trạng thái học sinh. & Bố trí danh sách, bộ lọc, nút thao tác và cảnh báo nhất quán. & \\
\hline
3 & Lớp và buổi học & Tạo lớp, lịch, buổi học, điểm danh. & Duy trì quy trình tuần tự từ lớp đến buổi học và điểm danh. & \\
\hline
4 & Tài chính & Nhập giao dịch, kiểm tra phí, xem số dư. & Tách rõ giao dịch thu/hoàn, khoản phí và báo cáo. & \\
\hline
5 & Tích hợp & Thực hiện OAuth, chọn guild/channel, gắn Gym. & Hiển thị trạng thái kết nối và lỗi từ dịch vụ ngoài dễ hiểu. & \\
\hline
\end{tabularx}
\end{table}

\subsubsection{Yêu cầu tương thích}

Yêu cầu tương thích mô tả khả năng hoạt động của phần mềm trên các môi trường
trình duyệt, hệ điều hành và dịch vụ ngoài mà hệ thống cần kết nối.

\noindent - Bảng yêu cầu tương thích:

\begin{table}[!htbp]
\centering
\tablesetup
\begin{tabularx}{\textwidth}{|C{0.9cm}|L{3.1cm}|Y|L{1.8cm}|}
\hline
\textbf{STT} & \textbf{Nghiệp vụ} & \textbf{Đối tượng liên quan} & \textbf{Ghi chú} \\
\hline
1 & Ứng dụng web & Trình duyệt hiện đại như Chrome, Edge, Firefox; giao diện desktop/laptop và màn hình phổ biến. & React + Vite. \\
\hline
2 & Backend API & Node.js, Express, TypeORM và Microsoft SQL Server/\allowbreak Azure SQL. & Modular monolith. \\
\hline
3 & Discord & Discord OAuth, bot token, guild, text channel, voice channel và DM. & Phụ thuộc API ngoài. \\
\hline
4 & Codeforces & Codeforces credential, Gym, \code{GymProblem} và standing. & Phụ thuộc API ngoài. \\
\hline
\end{tabularx}
\end{table}

\subsubsection{Yêu cầu tiến hóa}

Yêu cầu tiến hóa mô tả khả năng thay đổi các quy định, tham số nghiệp vụ và cấu
hình vận hành mà không cần viết lại toàn bộ hệ thống.

\noindent - Bảng yêu cầu tiến hóa:

\begin{table}[!htbp]
\centering
\tablesetup
\begin{tabularx}{\textwidth}{|C{0.9cm}|L{3.4cm}|L{3.2cm}|Y|}
\hline
\textbf{STT} & \textbf{Nghiệp vụ} & \textbf{Tham số cần thay đổi} & \textbf{Miền giá trị cần thay đổi} \\
\hline
1 & Thay đổi học phí mỗi buổi của lớp & Học phí mỗi buổi & Số nguyên không âm theo từng lớp \\
\hline
2 & Thay đổi lịch học định kỳ & Thứ trong tuần, giờ bắt đầu, giờ kết thúc & Lịch hợp lệ theo tuần \\
\hline
3 & Thay đổi trạng thái học sinh & Quy trình đang học, chờ lưu trữ, lưu trữ & Trạng thái Enrollment theo công nợ \\
\hline
4 & Thay đổi trạng thái điểm danh & Danh mục trạng thái điểm danh & Có mặt, vắng có phép, vắng không phép hoặc mở rộng sau này \\
\hline
5 & Thay đổi chính sách tài chính & Loại giao dịch, trạng thái khoản phí, cách tính số dư & Thu tiền, hoàn tiền, phí hiệu lực/hủy \\
\hline
6 & Thay đổi cấu hình Discord & Bot credential, guild, channel, permissions, scopes & Theo từng môi trường triển khai \\
\hline
7 & Thay đổi chu kỳ đồng bộ Codeforces & Pull interval của Gym & Số nguyên dương theo Gym hoặc lớp \\
\hline
8 & Thay đổi kỳ báo cáo & Từ ngày, đến ngày, danh sách lớp, bao gồm khoản chưa thu & Khoảng thời gian và bộ lọc báo cáo \\
\hline
\end{tabularx}
\end{table}

\noindent - Bảng trách nhiệm yêu cầu tiến hóa:

\begin{table}[!htbp]
\centering
\scriptsize
\setlength{\tabcolsep}{2pt}
\renewcommand{\arraystretch}{1.2}
\begin{tabularx}{\textwidth}{|C{0.8cm}|L{2.7cm}|L{3.2cm}|Y|L{1.6cm}|}
\hline
\textbf{STT} & \textbf{Nghiệp vụ} & \textbf{Người dùng} & \textbf{Phần mềm} & \textbf{Ghi chú} \\
\hline
1 & Học phí và lịch học & Giáo viên cập nhật lớp và lịch. & Ghi nhận giá trị mới, áp dụng cho buổi học và khoản phí phát sinh sau đó. & \\
\hline
2 & Điểm danh và tài chính & Giáo viên điều chỉnh trạng thái điểm danh hoặc khoản phí. & Cập nhật số dư, \code{FeeRecord} và báo cáo dựa trên trạng thái mới. & \\
\hline
3 & Discord & Quản trị viên hoặc giáo viên cập nhật cấu hình liên quan. & Lưu credential, guild, channel và dùng cho các thao tác gửi tin về sau. & \\
\hline
4 & Codeforces & Giáo viên thay đổi Gym hoặc chu kỳ đồng bộ. & Cập nhật liên kết Gym và cơ chế đồng bộ nền. & \\
\hline
5 & Báo cáo & Giáo viên chọn kỳ hoặc phạm vi lớp khác. & Tính toán lại số liệu theo tham số lọc mới. & \\
\hline
\end{tabularx}
\end{table}

\subsection{Yêu cầu hệ thống}

\subsubsection{Yêu cầu bảo mật}

Yêu cầu bảo mật đảm bảo tính toàn vẹn và bảo mật của dữ liệu học sinh, lớp học,
điểm danh, tài chính, credential Discord/\allowbreak Codeforces và thiết lập cơ chế phân
quyền truy cập rõ ràng giữa giáo viên, quản trị viên và học sinh.

\noindent - Bảng yêu cầu bảo mật:

\begin{table}[!htbp]
\centering
\scriptsize
\setlength{\tabcolsep}{2pt}
\renewcommand{\arraystretch}{1.15}
\begin{tabularx}{\textwidth}{|C{0.7cm}|L{2.8cm}|Y|C{1cm}|C{1cm}|C{1cm}|L{1.4cm}|}
\hline
\textbf{STT} & \textbf{Nghiệp vụ} & \textbf{Thao tác cụ thể} & \textbf{Giáo viên} & \textbf{Admin} & \textbf{Học sinh} & \textbf{Ghi chú} \\
\hline
1 & Quản lý tài khoản & Đăng nhập, xem/cập nhật hồ sơ cá nhân & X & X &  & Theo JWT và vai trò \\
\hline
2 & Quản trị tài khoản giáo viên & Xem danh sách, cập nhật, vô hiệu hóa giáo viên &  & X &  & Chỉ admin \\
\hline
3 & Quản lý học sinh & Thêm/sửa/xem/chuyển/rút/lưu trữ học sinh & X &  &  & Theo ownership \\
\hline
4 & Quản lý lớp học & Thêm/sửa/lưu trữ lớp, lịch, buổi học & X &  &  & Theo ownership \\
\hline
5 & Điểm danh & Xem và cập nhật điểm danh & X &  &  & Theo session/student ownership \\
\hline
6 & Tài chính & Quản lý giao dịch, khoản phí, số dư, báo cáo & X &  &  & Dữ liệu nhạy cảm \\
\hline
7 & Discord giáo viên/lớp & Liên kết Discord, chọn guild/channel, gửi thông báo/DM & X &  &  & Phụ thuộc credential \\
\hline
8 & Cấu hình bot Discord & Cập nhật token, client id, client secret &  & X &  & Dữ liệu rất nhạy cảm \\
\hline
9 & Uỷ quyền Discord của học sinh & Học sinh xác nhận OAuth để lưu Discord identity/token &  &  & X & Học sinh là tác nhân ngoài \\
\hline
10 & Codeforces Gym & Cấu hình credential, gắn Gym, xem standing & X & X &  & Admin có thể cập nhật credential giáo viên \\
\hline
\end{tabularx}
\end{table}

\noindent - Bảng trách nhiệm yêu cầu bảo mật:

\begin{table}[!htbp]
\centering
\scriptsize
\setlength{\tabcolsep}{2pt}
\renewcommand{\arraystretch}{1.2}
\begin{tabularx}{\textwidth}{|C{0.8cm}|L{2.4cm}|L{3.5cm}|Y|L{1.5cm}|}
\hline
\textbf{STT} & \textbf{Vai trò} & \textbf{Người dùng} & \textbf{Phần mềm} & \textbf{Ghi chú} \\
\hline
1 & Giáo viên & Sử dụng dữ liệu học sinh, lớp, tài chính, Discord và Codeforces đúng phạm vi mình quản lý. & Bắt buộc xác thực JWT, kiểm tra vai trò giáo viên và kiểm tra ownership ở các thao tác theo lớp/học sinh/buổi học/giao dịch. & \\
\hline
2 & Quản trị viên & Quản lý tài khoản giáo viên và cấu hình bot Discord. & Bắt buộc vai trò admin, không cho truy cập nghiệp vụ giáo viên nếu không đúng quyền. & \\
\hline
3 & Hệ thống ngoài & Discord và Codeforces trả dữ liệu qua API/\allowbreak OAuth. & Kiểm tra callback, bảo vệ token/secret, ghi nhận lỗi tích hợp và không lộ credential ở phản hồi. & \\
\hline
\end{tabularx}
\end{table}

\subsubsection{Yêu cầu an toàn}

Yêu cầu an toàn tập trung vào việc bảo toàn dữ liệu nghiệp vụ quan trọng, tránh
xóa nhầm hoặc làm mất lịch sử học tập, điểm danh, tài chính và dữ liệu tích hợp
ngoài hệ thống.

\noindent - Bảng yêu cầu an toàn:

\begin{table}[!htbp]
\centering
\scriptsize
\setlength{\tabcolsep}{2pt}
\renewcommand{\arraystretch}{1.2}
\begin{tabularx}{\textwidth}{|C{0.8cm}|L{3.5cm}|L{3.4cm}|Y|}
\hline
\textbf{STT} & \textbf{Nghiệp vụ không được phép thực hiện} & \textbf{Đối tượng liên quan} & \textbf{Ghi chú} \\
\hline
1 & Xóa cứng tài khoản giáo viên đang có dữ liệu nghiệp vụ & Giáo viên, học sinh, lớp, tài chính & Chỉ nên vô hiệu hóa tài khoản để giữ lịch sử. \\
\hline
2 & Xóa cứng học sinh đã có ghi danh, điểm danh, khoản phí, giao dịch hoặc kết quả Gym & Học sinh & Dùng trạng thái rút/chờ lưu trữ/lưu trữ thay cho xóa. \\
\hline
3 & Xóa cứng lớp học đã có học sinh, buổi học, điểm danh, Discord binding hoặc Gym & Lớp học & Dùng trạng thái archived. \\
\hline
4 & Xóa buổi học đã có điểm danh hoặc khoản phí phát sinh & \code{Session}, \code{Attendance}, \code{FeeRecord} & Dùng trạng thái cancelled và xử lý khoản phí liên quan. \\
\hline
5 & Ghi giao dịch hoặc cập nhật khoản phí cho học sinh không thuộc giáo viên hiện tại & \code{Student}, \code{Transaction}, \code{FeeRecord} & Phải kiểm tra ownership. \\
\hline
6 & Gửi thông báo Discord đến guild/channel không thuộc cấu hình lớp của giáo viên & Discord guild/channel & Tránh gửi sai nơi nhận. \\
\hline
7 & Hiển thị đầy đủ bot token, client secret, Codeforces API key/secret & Credential & Chỉ lưu bảo mật, hiển thị dạng ẩn hoặc không hiển thị. \\
\hline
8 & Gắn Gym hoặc xem standing của lớp không thuộc giáo viên hiện tại & \code{Class}, Gym, \code{GymStanding} & Phải kiểm tra ownership. \\
\hline
\end{tabularx}
\end{table}

\noindent - Bảng trách nhiệm yêu cầu an toàn:

\begin{table}[!htbp]
\centering
\scriptsize
\setlength{\tabcolsep}{2pt}
\renewcommand{\arraystretch}{1.2}
\begin{tabularx}{\textwidth}{|C{0.8cm}|L{3cm}|L{3.2cm}|Y|L{1.5cm}|}
\hline
\textbf{STT} & \textbf{Nghiệp vụ} & \textbf{Người thực hiện} & \textbf{Biện pháp nghiệp vụ} & \textbf{Ghi chú} \\
\hline
1 & Kiểm soát tài khoản & Quản trị viên & Vô hiệu hóa thay vì xóa cứng, bảo vệ credential nhạy cảm. & \\
\hline
2 & Kiểm soát học sinh & Giáo viên & Quản lý vòng đời qua trạng thái ghi danh và lưu trữ. & \\
\hline
3 & Kiểm soát lớp và buổi học & Giáo viên & Lưu trữ lớp, hủy buổi học bằng trạng thái, giữ dữ liệu lịch sử. & \\
\hline
4 & Kiểm soát tài chính & Giáo viên & Không xóa mất dấu vết phí/giao dịch; cập nhật trạng thái thay cho mất lịch sử. & \\
\hline
5 & Kiểm soát tích hợp ngoài & Giáo viên, quản trị viên & Kiểm tra quyền sở hữu guild, channel, Gym và xử lý lỗi API ngoài. & \\
\hline
\end{tabularx}
\end{table}

\subsection{Yêu cầu công nghệ}

\noindent - Bảng yêu cầu công nghệ:

\begin{table}[!htbp]
\centering
\scriptsize
\setlength{\tabcolsep}{2pt}
\renewcommand{\arraystretch}{1.15}
\begin{tabularx}{\textwidth}{|C{0.9cm}|L{3cm}|Y|L{2.4cm}|}
\hline
\textbf{STT} & \textbf{Yêu cầu} & \textbf{Mô tả chi tiết} & \textbf{Ghi chú} \\
\hline
1 & Nền tảng & Ứng dụng web chạy trên trình duyệt hiện đại. & Desktop/laptop là chính. \\
\hline
2 & Frontend & React với Vite, TypeScript, component UI dùng trong thư mục frontend. & Giao diện tương tác. \\
\hline
3 & Backend & Node.js, Express, TypeScript, TypeORM theo kiến trúc modular monolith. & API dưới tiền tố \code{/api}. \\
\hline
4 & Cơ sở dữ liệu & Microsoft SQL Server/\allowbreak Azure SQL qua driver \code{mssql}, lưu Teacher, \code{Student}, \code{Class}, \code{Session}, \code{Attendance}, \code{Transaction}, \code{FeeRecord}, Discord và Gym. & Dữ liệu quan hệ. \\
\hline
5 & Bảo mật & JWT, bcrypt, RBAC, kiểm tra ownership, mã hóa hoặc che giấu credential nhạy cảm. & Tài khoản và credential. \\
\hline
6 & Tích hợp Discord & Discord OAuth, bot credential, guild/channel cache, gửi channel post và DM. & API ngoài. \\
\hline
7 & Tích hợp Codeforces & Codeforces credential, đồng bộ Gym, \code{GymProblem}, \code{GymStanding}. & API ngoài. \\
\hline
8 & Triển khai & Có Dockerfile và docker-compose để đóng gói ứng dụng; cơ sở dữ liệu được cấu hình qua biến môi trường DB hoặc chuỗi kết nối Azure SQL. & Phù hợp máy chủ nội bộ hoặc cloud. \\
\hline
9 & Dễ bảo trì & Phân module Account, System, \code{Student}, Classroom, Finance; backend giao tiếp nội bộ qua contract/persistence thay vì HTTP nội bộ. & Giảm phụ thuộc chéo. \\
\hline
\end{tabularx}
\end{table}

\noindent - Bảng trách nhiệm yêu cầu công nghệ:

\begin{table}[!htbp]
\centering
\scriptsize
\setlength{\tabcolsep}{2pt}
\renewcommand{\arraystretch}{1.2}
\begin{tabularx}{\textwidth}{|C{0.9cm}|L{3cm}|L{3.4cm}|Y|L{1.7cm}|}
\hline
\textbf{STT} & \textbf{Nghiệp vụ} & \textbf{Người dùng} & \textbf{Phần mềm} & \textbf{Ghi chú} \\
\hline
1 & Vận hành frontend & Giáo viên, quản trị viên & Cung cấp giao diện web ổn định trên trình duyệt hiện đại. & React/\allowbreak Vite. \\
\hline
2 & Vận hành backend &  & Cung cấp API ổn định, phân quyền, kiểm tra dữ liệu và xử lý nghiệp vụ. & Express/\allowbreak TypeORM. \\
\hline
3 & Vận hành cơ sở dữ liệu &  & Lưu trữ và truy xuất chính xác dữ liệu học sinh, lớp, buổi học, điểm danh, tài chính và tích hợp. & Azure SQL/MSSQL. \\
\hline
4 & Vận hành tích hợp & Giáo viên, quản trị viên & Kết nối Discord và Codeforces, cache dữ liệu cần thiết và ghi nhận lỗi. & API ngoài. \\
\hline
\end{tabularx}
\end{table}

\subsection{Phân tích yêu cầu}

\begin{itemize}[leftmargin=*]
\item \textbf{Tính khả thi kỹ thuật:} Khá khả thi. Dự án hiện tại đã có cấu trúc backend theo modular monolith với các module Account, System, \code{Student}, Classroom và Finance; frontend dùng React/\allowbreak Vite; cơ sở dữ liệu dùng Microsoft SQL Server/\allowbreak Azure SQL thông qua TypeORM và driver \code{mssql}. Các nghiệp vụ nội bộ như quản lý học sinh, ghi danh, lớp học, buổi học, điểm danh và tài chính là các chức năng web quản trị phổ biến. Rủi ro kỹ thuật chính nằm ở tích hợp Discord và Codeforces vì phụ thuộc API ngoài, credential, OAuth, giới hạn tốc độ và khả năng lỗi mạng.

\item \textbf{Tính khả thi nghiệp vụ:} Hoàn toàn khả thi. Quy trình nghiệp vụ đã rõ hơn so với bản trước: giáo viên quản lý học sinh theo Enrollment, tổ chức lớp bằng \code{ClassSchedule} và \code{Session}, ghi nhận \code{Attendance}, phát sinh \code{FeeRecord}, đối soát bằng \code{Transaction} và xem báo cáo. Codeforces không còn chỉ là thông tin rời rạc mà được mô hình hóa thành Gym, \code{GymProblem} và \code{GymStanding}; Discord không chỉ gửi tin mà còn có bot credential, guild binding, channel thông báo và voice channel điểm danh.

\item \textbf{Tính khả thi về con người:} Khá khả thi nhưng cần hướng dẫn ban đầu cho các luồng nhiều bước. Giáo viên có thể làm quen nhanh với các chức năng cốt lõi nếu giao diện đi theo quy trình lớp học thực tế: tạo lớp, thêm học sinh, tạo buổi học, điểm danh, xem phí và thu tiền. Các phần khó hơn là cấu hình Discord, liên kết OAuth và Codeforces credential, nên cần thông báo lỗi rõ ràng và tách quyền quản trị viên khỏi quyền giáo viên.

\item \textbf{Rủi ro dữ liệu:} Dữ liệu tài chính, điểm danh và trạng thái ghi danh có quan hệ chặt chẽ. Vì vậy hệ thống nên ưu tiên cập nhật bằng trạng thái thay vì xóa cứng, kiểm tra ownership ở mọi thao tác có khóa ngoại và giữ lịch sử đủ để truy vết học phí, hoàn tiền, buổi học đã hủy hoặc học sinh đã lưu trữ.

\item \textbf{Rủi ro mở rộng:} Khi số lượng buổi học, điểm danh, \code{FeeRecord} và \code{GymStanding} tăng, các màn hình báo cáo cần được phân trang, lọc theo thời gian/lớp và tối ưu truy vấn. Những tác vụ đồng bộ Codeforces hoặc Discord nên được xử lý nền hoặc dùng cache để giao diện không phụ thuộc hoàn toàn vào thời gian phản hồi của API ngoài.
\end{itemize}

\setcounter{secnumdepth}{5}
\renewcommand{\theparagraph}{\thesubsubsection.\arabic{paragraph}}
\renewcommand{\thesubparagraph}{\theparagraph.\arabic{subparagraph}}
\titleformat{\paragraph}[block]{\normalsize\itshape\bfseries}{\theparagraph.}{0.5em}{}
\titlespacing*{\paragraph}{0pt}{6pt}{0pt}
\titleformat{\subparagraph}[block]{\normalsize\itshape\bfseries}{\thesubparagraph.}{0.5em}{}
\titlespacing*{\subparagraph}{0pt}{6pt}{0pt}

\section{Mô hình hoá yêu cầu chức năng}

\subsection{Tổng quan về mô hình hoá yêu cầu}

Mục III là tiền đề trực tiếp để xây dựng mục IV. Các biểu mẫu BM 1.1 đến BM
1.25, các quy định QĐ 1.1 đến QĐ 1.25 và bảng trách nhiệm yêu cầu nghiệp vụ ở
mục III được dùng để xác định tác nhân, dữ liệu nhập, dữ liệu xuất, kho dữ liệu,
điều kiện kiểm tra và các nhánh lỗi của từng mô hình. Vì vậy, phần mô hình hoá
không mô tả lại một hệ thống cũ mà bám theo dự án hiện tại: hệ thống giúp giáo
viên quản lý lớp học lập trình thi đấu.

Hệ thống có năm nhóm tác nhân chính:

\begin{table}[!htbp]
\centering
\tablesetup
\begin{tabularx}{\textwidth}{|C{0.8cm}|L{2.6cm}|Y|}
\hline
\textbf{STT} & \textbf{Tác nhân} & \textbf{Vai trò} \\
\hline
1 & Giáo viên & Người dùng chính, quản lý học sinh, ghi danh, lớp học, lịch học, buổi học, điểm danh, tài chính, Discord và Codeforces Gym. \\
\hline
2 & Quản trị viên & Quản lý tài khoản giáo viên và cấu hình bot Discord toàn hệ thống. \\
\hline
3 & Học sinh & Thực hiện uỷ quyền Discord OAuth để hệ thống lưu \code{StudentDiscordCredential}. \\
\hline
4 & Discord & Hệ thống ngoài cung cấp OAuth, guild, channel, member, voice state và API gửi tin. \\
\hline
5 & Codeforces & Hệ thống ngoài cung cấp dữ liệu Gym, bài trong Gym và bảng xếp hạng Gym. \\
\hline
\end{tabularx}
\end{table}

\subsection{Mô hình hoá theo chức năng (DFD -- Data Flow Diagram)}

Các sơ đồ DFD dưới đây dùng cùng một quy ước: D1 là dữ liệu/yêu cầu từ tác
nhân, D2 là dữ liệu từ thiết bị nhập hoặc hệ thống ngoài; nếu không phát sinh
dữ liệu riêng ngoài D1 thì ghi ``Không có'', D3 là dữ liệu đọc từ bộ nhớ phụ,
D4 là dữ liệu ghi vào bộ nhớ phụ, D5 là dữ liệu hiển thị hoặc kết xuất; nếu
không có dữ liệu xuất riêng thì ghi ``Không có'', D6 là thông báo kết quả trả
về tác nhân. ``Bộ nhớ phụ'' trong các sơ đồ tương
ứng với cơ sở dữ liệu Microsoft SQL Server/\allowbreak Azure SQL của hệ thống.

\subsubsection{Quản lý tài khoản}

\paragraph{Chức năng Đăng nhập hệ thống}
\dfddiagram{Đăng nhập\\hệ thống}{Giáo viên/\\Quản trị viên}{Thiết bị nhập}{Teacher/Admin}
\textbf{Mô tả luồng dữ liệu:}
\begin{itemize}[leftmargin=*]
\item \textbf{D1:} Tên đăng nhập và mật khẩu theo BM 1.1.
\item \textbf{D2:} Không có.
\item \textbf{D3:} Tài khoản, mật khẩu đã băm, vai trò và trạng thái hoạt động.
\item \textbf{D4:} Không có.
\item \textbf{D5:} JWT, vai trò và dữ liệu điều hướng.
\item \textbf{D6:} Thông báo đăng nhập thành công hoặc lỗi theo QĐ 1.1.
\end{itemize}
\textbf{Thuật toán:}
\begin{itemize}[leftmargin=*]
\item \textbf{B1:} Nhận tên đăng nhập và mật khẩu.
\item \textbf{B2:} Tìm tài khoản tương ứng trong bộ nhớ phụ.
\item \textbf{B3:} Kiểm tra mật khẩu, trạng thái hoạt động và vai trò.
\item \textbf{B4:} Nếu không hợp lệ, trả lỗi và kết thúc.
\item \textbf{B5:} Nếu hợp lệ, cấp token và điều hướng vào màn hình phù hợp.
\end{itemize}

\paragraph{Chức năng Đăng ký tài khoản giáo viên}
\dfddiagram{Đăng ký\\tài khoản giáo viên}{Quản trị viên/\\Giáo viên}{Thiết bị nhập}{Teacher}
\textbf{Mô tả luồng dữ liệu:}
\begin{itemize}[leftmargin=*]
\item \textbf{D1:} Tên đăng nhập, mật khẩu, Codeforces handle/API key/API secret theo BM 1.2.
\item \textbf{D2:} Không có.
\item \textbf{D3:} Danh sách tài khoản giáo viên hiện có.
\item \textbf{D4:} Tài khoản giáo viên mới và credential Codeforces được lưu bảo mật.
\item \textbf{D5:} Giao diện danh sách tài khoản hoặc màn hình đăng nhập.
\item \textbf{D6:} Thông báo đăng ký thành công hoặc lỗi theo QĐ 1.2.
\end{itemize}
\textbf{Thuật toán:}
\begin{itemize}[leftmargin=*]
\item \textbf{B1:} Nhận dữ liệu đăng ký từ biểu mẫu.
\item \textbf{B2:} Kiểm tra tên đăng nhập và mật khẩu không rỗng.
\item \textbf{B3:} Kiểm tra tên đăng nhập không trùng.
\item \textbf{B4:} Băm mật khẩu, bảo vệ credential Codeforces nếu có.
\item \textbf{B5:} Lưu tài khoản và trả kết quả.
\end{itemize}

\paragraph{Chức năng Cập nhật hồ sơ và thông tin Codeforces của giáo viên}
\dfddiagram{Cập nhật hồ sơ\\giáo viên}{Giáo viên/\\Quản trị viên}{Thiết bị nhập}{Teacher, Credential}
\textbf{Mô tả luồng dữ liệu:}
\begin{itemize}[leftmargin=*]
\item \textbf{D1:} Các trường cập nhật theo BM 1.3.
\item \textbf{D2:} Không có.
\item \textbf{D3:} Hồ sơ giáo viên và credential Codeforces hiện có.
\item \textbf{D4:} Hồ sơ hoặc credential mới sau khi cập nhật.
\item \textbf{D5:} Hồ sơ đã được che thông tin nhạy cảm.
\item \textbf{D6:} Thông báo thành công hoặc lỗi theo QĐ 1.3.
\end{itemize}
\textbf{Thuật toán:}
\begin{itemize}[leftmargin=*]
\item \textbf{B1:} Nhận yêu cầu cập nhật hồ sơ.
\item \textbf{B2:} Kiểm tra có ít nhất một trường thay đổi.
\item \textbf{B3:} Nếu cập nhật mật khẩu thì mật khẩu không được rỗng.
\item \textbf{B4:} Lưu dữ liệu mới, che API key/secret khi trả về giao diện.
\end{itemize}

\paragraph{Chức năng Liên kết tài khoản Discord của giáo viên}
\dfddiagram{Liên kết Discord\\giáo viên}{Giáo viên}{Discord OAuth}{Teacher Discord identity}
\textbf{Mô tả luồng dữ liệu:}
\begin{itemize}[leftmargin=*]
\item \textbf{D1:} Yêu cầu xác minh Discord theo BM 1.4.
\item \textbf{D2:} OAuth callback và thông tin user từ Discord.
\item \textbf{D3:} Hồ sơ giáo viên hiện tại và cấu hình OAuth.
\item \textbf{D4:} Discord user id, username của giáo viên.
\item \textbf{D5:} Trạng thái xác minh Discord.
\item \textbf{D6:} Kết quả liên kết hoặc lỗi callback theo QĐ 1.4.
\end{itemize}
\textbf{Thuật toán:}
\begin{itemize}[leftmargin=*]
\item \textbf{B1:} Giáo viên yêu cầu liên kết Discord.
\item \textbf{B2:} Hệ thống sinh URL OAuth.
\item \textbf{B3:} Discord trả callback.
\item \textbf{B4:} Hệ thống xác thực callback và lưu Discord identity.
\end{itemize}

\subsubsection{Quản lý học sinh và ghi danh}

\paragraph{Chức năng Thêm học sinh và ghi danh vào lớp}
\dfddiagram{Thêm học sinh\\và ghi danh}{Giáo viên}{Thiết bị nhập}{Student, Enrollment, Class}
\textbf{Mô tả luồng dữ liệu:}
\begin{itemize}[leftmargin=*]
\item \textbf{D1:} Họ tên, lớp học, Codeforces handle, số điện thoại, ghi chú theo BM 1.5.
\item \textbf{D2:} Không có.
\item \textbf{D3:} Lớp học thuộc giáo viên, danh sách học sinh và ghi danh hiện có.
\item \textbf{D4:} Hồ sơ Student và Enrollment mới.
\item \textbf{D5:} Danh sách học sinh sau khi thêm.
\item \textbf{D6:} Thông báo kết quả theo QĐ 1.5.
\end{itemize}
\textbf{Thuật toán:}
\begin{itemize}[leftmargin=*]
\item \textbf{B1:} Nhận dữ liệu học sinh và lớp ghi danh.
\item \textbf{B2:} Kiểm tra lớp thuộc giáo viên hiện tại.
\item \textbf{B3:} Kiểm tra họ tên, lớp học, Codeforces handle không rỗng.
\item \textbf{B4:} Kiểm tra học sinh không có ghi danh hoạt động bị trùng trong phạm vi giáo viên.
\item \textbf{B5:} Tạo Student và Enrollment.
\end{itemize}

\paragraph{Chức năng Cập nhật hồ sơ học sinh}
\dfddiagram{Cập nhật\\hồ sơ học sinh}{Giáo viên}{Thiết bị nhập}{Student}
\textbf{Mô tả luồng dữ liệu:}
\begin{itemize}[leftmargin=*]
\item \textbf{D1:} Thông tin cập nhật theo BM 1.6.
\item \textbf{D2:} Không có.
\item \textbf{D3:} Hồ sơ học sinh thuộc giáo viên hiện tại.
\item \textbf{D4:} Hồ sơ học sinh đã cập nhật.
\item \textbf{D5:} Hồ sơ hoặc danh sách học sinh mới.
\item \textbf{D6:} Thông báo kết quả theo QĐ 1.6.
\end{itemize}
\textbf{Thuật toán:}
\begin{itemize}[leftmargin=*]
\item \textbf{B1:} Nhận yêu cầu cập nhật.
\item \textbf{B2:} Kiểm tra học sinh thuộc giáo viên hiện tại.
\item \textbf{B3:} Kiểm tra có ít nhất một trường thay đổi.
\item \textbf{B4:} Nếu cập nhật Codeforces handle thì handle không được rỗng.
\item \textbf{B5:} Lưu hồ sơ.
\end{itemize}

\paragraph{Chức năng Chuyển lớp, rút học sinh, khôi phục học sinh}
\dfddiagram{Chuyển lớp, rút,\\khôi phục học sinh}{Giáo viên}{Thiết bị nhập}{Student, Enrollment, Balance}
\textbf{Mô tả luồng dữ liệu:}
\begin{itemize}[leftmargin=*]
\item \textbf{D1:} Học sinh, lớp chuyển đến và thời điểm thao tác theo BM 1.7.
\item \textbf{D2:} Không có.
\item \textbf{D3:} Enrollment hiện tại, lớp đích và số dư học sinh.
\item \textbf{D4:} Enrollment cũ được kết thúc, Enrollment mới hoặc trạng thái rút/khôi phục.
\item \textbf{D5:} Trạng thái ghi danh mới.
\item \textbf{D6:} Thông báo thành công hoặc lỗi theo QĐ 1.7.
\end{itemize}
\textbf{Thuật toán:}
\begin{itemize}[leftmargin=*]
\item \textbf{B1:} Nhận thao tác chuyển/rút/khôi phục.
\item \textbf{B2:} Kiểm tra học sinh và lớp thuộc giáo viên.
\item \textbf{B3:} Nếu chuyển lớp, kết thúc ghi danh cũ và tạo ghi danh mới.
\item \textbf{B4:} Nếu rút học sinh, xác định trạng thái tiếp theo theo số dư.
\item \textbf{B5:} Lưu thay đổi và trả trạng thái mới.
\end{itemize}

\paragraph{Chức năng Lưu trữ học sinh đang chờ xử lý công nợ}
\dfddiagram{Lưu trữ học sinh\\chờ xử lý}{Giáo viên}{Thiết bị nhập}{Student, Transaction, FeeRecord}
\textbf{Mô tả luồng dữ liệu:}
\begin{itemize}[leftmargin=*]
\item \textbf{D1:} Học sinh, thời điểm lưu trữ, lý do chờ lưu trữ theo BM 1.8.
\item \textbf{D2:} Không có.
\item \textbf{D3:} Trạng thái học sinh và số dư hiện tại.
\item \textbf{D4:} Trạng thái học sinh đã lưu trữ hoặc tiếp tục chờ xử lý.
\item \textbf{D5:} Danh sách học sinh theo trạng thái.
\item \textbf{D6:} Kết quả theo QĐ 1.8.
\end{itemize}
\textbf{Thuật toán:}
\begin{itemize}[leftmargin=*]
\item \textbf{B1:} Nhận yêu cầu lưu trữ.
\item \textbf{B2:} Kiểm tra học sinh đang ở trạng thái chờ lưu trữ.
\item \textbf{B3:} Đọc số dư cần thu hoặc cần hoàn.
\item \textbf{B4:} Nếu chưa xử lý xong số dư, báo lý do.
\item \textbf{B5:} Nếu đủ điều kiện, chuyển trạng thái đã lưu trữ.
\end{itemize}

\paragraph{Chức năng Tra cứu và xem hồ sơ học tập của học sinh}
\dfddiagram{Tra cứu và xem\\hồ sơ học tập}{Giáo viên}{Thiết bị nhập}{Student, Enrollment, GymStanding}
\textbf{Mô tả luồng dữ liệu:}
\begin{itemize}[leftmargin=*]
\item \textbf{D1:} Từ khóa, lớp, trạng thái học sinh hoặc Codeforces handle.
\item \textbf{D2:} Không có.
\item \textbf{D3:} Hồ sơ học sinh, ghi danh, điểm danh, tài chính, GymStanding.
\item \textbf{D4:} Không có.
\item \textbf{D5:} Hồ sơ học tập tổng hợp.
\item \textbf{D6:} Kết quả tra cứu hoặc thông báo không có dữ liệu.
\end{itemize}
\textbf{Thuật toán:}
\begin{itemize}[leftmargin=*]
\item \textbf{B1:} Nhận bộ lọc.
\item \textbf{B2:} Truy vấn dữ liệu trong phạm vi giáo viên.
\item \textbf{B3:} Tổng hợp hồ sơ, ghi danh, điểm danh, công nợ và kết quả Gym.
\item \textbf{B4:} Trả danh sách hoặc hồ sơ chi tiết.
\end{itemize}

\subsubsection{Quản lý lớp học, lịch học và buổi học}

\paragraph{Chức năng Thêm/sửa/lưu trữ lớp học}
\dfddiagram{Thêm/sửa/lưu trữ\\lớp học}{Giáo viên}{Thiết bị nhập}{Class, ClassSchedule}
\textbf{Mô tả luồng dữ liệu:}
\begin{itemize}[leftmargin=*]
\item \textbf{D1:} Tên lớp, học phí mỗi buổi, lịch học, trạng thái lớp theo BM 1.10.
\item \textbf{D2:} Không có.
\item \textbf{D3:} Danh sách lớp và lịch học hiện có.
\item \textbf{D4:} Class và ClassSchedule mới/cập nhật hoặc trạng thái archived.
\item \textbf{D5:} Danh sách lớp học.
\item \textbf{D6:} Kết quả theo QĐ 1.10.
\end{itemize}
\textbf{Thuật toán:}
\begin{itemize}[leftmargin=*]
\item \textbf{B1:} Nhận thông tin lớp.
\item \textbf{B2:} Kiểm tra tên lớp không rỗng, học phí không âm.
\item \textbf{B3:} Kiểm tra lớp mới có ít nhất một lịch học.
\item \textbf{B4:} Lưu lớp và lịch học; nếu lưu trữ thì chuyển trạng thái thay vì xóa.
\end{itemize}

\paragraph{Chức năng Thiết lập lịch học định kỳ của lớp}
\dfddiagram{Thiết lập\\lịch học định kỳ}{Giáo viên}{Thiết bị nhập}{ClassSchedule}
\textbf{Mô tả luồng dữ liệu:}
\begin{itemize}[leftmargin=*]
\item \textbf{D1:} Thứ trong tuần, giờ bắt đầu, giờ kết thúc theo BM 1.11.
\item \textbf{D2:} Không có.
\item \textbf{D3:} Lớp học và lịch học hiện có.
\item \textbf{D4:} Lịch học định kỳ mới thay thế cấu hình đang áp dụng.
\item \textbf{D5:} Lịch học của lớp.
\item \textbf{D6:} Kết quả theo QĐ 1.11.
\end{itemize}
\textbf{Thuật toán:}
\begin{itemize}[leftmargin=*]
\item \textbf{B1:} Nhận danh sách lịch.
\item \textbf{B2:} Kiểm tra thứ trong tuần hợp lệ.
\item \textbf{B3:} Kiểm tra giờ kết thúc sau giờ bắt đầu.
\item \textbf{B4:} Thay thế lịch đang áp dụng của lớp.
\end{itemize}

\paragraph{Chức năng Tạo buổi học thủ công và xem danh sách buổi học}
\dfddiagram{Tạo và xem\\buổi học}{Giáo viên}{Thiết bị nhập}{Session, Class}
\textbf{Mô tả luồng dữ liệu:}
\begin{itemize}[leftmargin=*]
\item \textbf{D1:} Lớp học, ngày học, giờ bắt đầu, giờ kết thúc theo BM 1.12.
\item \textbf{D2:} Không có.
\item \textbf{D3:} Lớp, lịch học, buổi học hiện có.
\item \textbf{D4:} Session mới cần lưu trữ; nếu chỉ xem danh sách thì không có dữ liệu ghi.
\item \textbf{D5:} Danh sách buổi học.
\item \textbf{D6:} Kết quả theo QĐ 1.12.
\end{itemize}
\textbf{Thuật toán:}
\begin{itemize}[leftmargin=*]
\item \textbf{B1:} Nhận yêu cầu tạo hoặc tra cứu buổi học.
\item \textbf{B2:} Kiểm tra lớp thuộc giáo viên.
\item \textbf{B3:} Kiểm tra đủ nhóm dữ liệu thời gian.
\item \textbf{B4:} Kiểm tra giờ kết thúc sau giờ bắt đầu.
\item \textbf{B5:} Tạo Session hoặc trả danh sách Session.
\end{itemize}

\paragraph{Chức năng Hủy buổi học}
\dfddiagram{Hủy\\buổi học}{Giáo viên}{Thiết bị nhập}{Session, Attendance, FeeRecord}
\textbf{Mô tả luồng dữ liệu:}
\begin{itemize}[leftmargin=*]
\item \textbf{D1:} Buổi học cần hủy.
\item \textbf{D2:} Không có.
\item \textbf{D3:} Buổi học, điểm danh và khoản phí liên quan.
\item \textbf{D4:} Session chuyển sang trạng thái hủy và dữ liệu liên quan được xử lý.
\item \textbf{D5:} Trạng thái buổi học sau khi hủy.
\item \textbf{D6:} Kết quả theo QĐ 1.13.
\end{itemize}
\textbf{Thuật toán:}
\begin{itemize}[leftmargin=*]
\item \textbf{B1:} Giáo viên chọn buổi học cần hủy.
\item \textbf{B2:} Kiểm tra buổi học thuộc giáo viên hiện tại.
\item \textbf{B3:} Chuyển trạng thái buổi học sang hủy, không xóa cứng.
\item \textbf{B4:} Xử lý Attendance và FeeRecord theo trạng thái buổi học.
\end{itemize}

\subsubsection{Quản lý điểm danh}

\paragraph{Chức năng Xem và cập nhật trạng thái điểm danh học sinh}
\dfddiagram{Cập nhật\\điểm danh}{Giáo viên/\\Discord}{Thiết bị nhập/\\Voice state}{Attendance, Session, FeeRecord}
\textbf{Mô tả luồng dữ liệu:}
\begin{itemize}[leftmargin=*]
\item \textbf{D1:} Buổi học, học sinh, trạng thái điểm danh, ghi chú theo BM 1.14.
\item \textbf{D2:} Voice state từ Discord; nếu giáo viên cập nhật thủ công thì không có.
\item \textbf{D3:} Session, Enrollment, Attendance, FeeRecord hiện có.
\item \textbf{D4:} Attendance cập nhật và FeeRecord phát sinh/hủy/kích hoạt lại.
\item \textbf{D5:} Bảng điểm danh và khoản phí liên quan.
\item \textbf{D6:} Kết quả theo QĐ 1.14.
\end{itemize}
\textbf{Thuật toán:}
\begin{itemize}[leftmargin=*]
\item \textbf{B1:} Nhận trạng thái điểm danh.
\item \textbf{B2:} Kiểm tra học sinh và buổi học thuộc giáo viên.
\item \textbf{B3:} Kiểm tra trạng thái nằm trong ba giá trị hợp lệ.
\item \textbf{B4:} Lưu Attendance.
\item \textbf{B5:} Cập nhật FeeRecord tương ứng nếu điểm danh ảnh hưởng học phí.
\end{itemize}

\subsubsection{Quản lý tài chính và công nợ}

\paragraph{Chức năng Ghi nhận giao dịch thu tiền hoặc hoàn tiền}
\dfddiagram{Ghi nhận\\giao dịch}{Giáo viên}{Thiết bị nhập}{Transaction, Student}
\textbf{Mô tả luồng dữ liệu:}
\begin{itemize}[leftmargin=*]
\item \textbf{D1:} Học sinh, số tiền, loại giao dịch, ghi chú, thời điểm theo BM 1.15.
\item \textbf{D2:} Không có.
\item \textbf{D3:} Học sinh, giao dịch cũ nếu cập nhật và số dư hiện tại.
\item \textbf{D4:} Transaction mới/cập nhật.
\item \textbf{D5:} Lịch sử giao dịch và số dư mới.
\item \textbf{D6:} Kết quả theo QĐ 1.15.
\end{itemize}
\textbf{Thuật toán:}
\begin{itemize}[leftmargin=*]
\item \textbf{B1:} Nhận dữ liệu giao dịch.
\item \textbf{B2:} Kiểm tra học sinh thuộc giáo viên.
\item \textbf{B3:} Kiểm tra loại giao dịch là thu tiền hoặc hoàn tiền.
\item \textbf{B4:} Lưu thu tiền là số dương, hoàn tiền là số âm.
\item \textbf{B5:} Cập nhật số dư hiển thị.
\end{itemize}

\paragraph{Chức năng Cập nhật trạng thái khoản phí phát sinh từ buổi học}
\dfddiagram{Cập nhật\\khoản phí}{Giáo viên}{Thiết bị nhập}{FeeRecord}
\textbf{Mô tả luồng dữ liệu:}
\begin{itemize}[leftmargin=*]
\item \textbf{D1:} Khoản phí và trạng thái khoản phí theo BM 1.16.
\item \textbf{D2:} Không có.
\item \textbf{D3:} FeeRecord thuộc học sinh/buổi học của giáo viên.
\item \textbf{D4:} FeeRecord đang hiệu lực hoặc đã hủy.
\item \textbf{D5:} Danh sách khoản phí và số dư.
\item \textbf{D6:} Kết quả theo QĐ 1.16.
\end{itemize}
\textbf{Thuật toán:}
\begin{itemize}[leftmargin=*]
\item \textbf{B1:} Nhận khoản phí cần cập nhật.
\item \textbf{B2:} Kiểm tra quyền truy cập khoản phí.
\item \textbf{B3:} Kiểm tra trạng thái chỉ là đang hiệu lực hoặc đã hủy.
\item \textbf{B4:} Lưu trạng thái mới, không xóa lịch sử.
\end{itemize}

\paragraph{Chức năng Xem số dư học sinh và tổng quan tài chính}
\dfddiagram{Xem số dư và\\tổng quan tài chính}{Giáo viên}{Thiết bị nhập}{Transaction, FeeRecord}
\textbf{Mô tả luồng dữ liệu:}
\begin{itemize}[leftmargin=*]
\item \textbf{D1:} Bộ lọc học sinh, lớp, khoảng thời gian.
\item \textbf{D2:} Không có.
\item \textbf{D3:} Transaction, FeeRecord, Student, Class.
\item \textbf{D4:} Không có.
\item \textbf{D5:} Số dư học sinh, tổng phải thu, tổng đã thu, tổng hoàn.
\item \textbf{D6:} Kết quả đối soát.
\end{itemize}
\textbf{Thuật toán:}
\begin{itemize}[leftmargin=*]
\item \textbf{B1:} Nhận bộ lọc.
\item \textbf{B2:} Truy vấn khoản phí và giao dịch trong phạm vi giáo viên.
\item \textbf{B3:} Tổng hợp số dư theo học sinh/lớp.
\item \textbf{B4:} Trả bảng tổng quan.
\end{itemize}

\paragraph{Chức năng Lập báo cáo thu nhập theo kỳ}
\dfddiagram{Lập báo cáo\\thu nhập}{Giáo viên}{Thiết bị nhập}{Transaction, FeeRecord, Class}
\textbf{Mô tả luồng dữ liệu:}
\begin{itemize}[leftmargin=*]
\item \textbf{D1:} Từ ngày, đến ngày, danh sách lớp, tùy chọn bao gồm khoản chưa thu theo BM 1.18.
\item \textbf{D2:} Không có.
\item \textbf{D3:} Dữ liệu giao dịch, khoản phí, lớp và học sinh.
\item \textbf{D4:} Không có.
\item \textbf{D5:} Báo cáo thu nhập theo kỳ.
\item \textbf{D6:} Kết quả theo QĐ 1.18.
\end{itemize}
\textbf{Thuật toán:}
\begin{itemize}[leftmargin=*]
\item \textbf{B1:} Nhận bộ lọc báo cáo.
\item \textbf{B2:} Kiểm tra từ ngày không sau đến ngày.
\item \textbf{B3:} Kiểm tra lớp thuộc giáo viên.
\item \textbf{B4:} Tổng hợp giao dịch, khoản phí, số dư.
\item \textbf{B5:} Trả báo cáo hoặc file kết xuất.
\end{itemize}

\subsubsection{Quản lý Discord}

\paragraph{Chức năng Cấu hình bot Discord toàn hệ thống}
\dfddiagram{Cấu hình\\bot Discord}{Quản trị viên}{Thiết bị nhập}{SysadminDiscordBotCredential}
\textbf{Mô tả luồng dữ liệu:}
\begin{itemize}[leftmargin=*]
\item \textbf{D1:} Bot token, client id, client secret, permissions, scopes theo BM 1.19.
\item \textbf{D2:} Không có.
\item \textbf{D3:} Cấu hình bot hiện có.
\item \textbf{D4:} Credential bot được lưu bảo mật.
\item \textbf{D5:} Thông tin cấu hình đã che dữ liệu nhạy cảm.
\item \textbf{D6:} Kết quả theo QĐ 1.19.
\end{itemize}
\textbf{Thuật toán:}
\begin{itemize}[leftmargin=*]
\item \textbf{B1:} Quản trị viên nhập cấu hình.
\item \textbf{B2:} Kiểm tra vai trò quản trị viên.
\item \textbf{B3:} Kiểm tra token, client id, client secret không rỗng.
\item \textbf{B4:} Lưu bảo mật và không trả giá trị đầy đủ ra giao diện.
\end{itemize}

\paragraph{Chức năng Chọn Discord guild và channel cho lớp}
\dfddiagram{Chọn guild/channel\\cho lớp}{Giáo viên}{Discord API}{ClassDiscordBinding}
\textbf{Mô tả luồng dữ liệu:}
\begin{itemize}[leftmargin=*]
\item \textbf{D1:} Lớp học, guild, channel thông báo, voice channel theo BM 1.20.
\item \textbf{D2:} Danh sách guild/channel từ Discord.
\item \textbf{D3:} Lớp, Discord guild/channel cache, binding hiện có.
\item \textbf{D4:} ClassDiscordBinding mới/cập nhật.
\item \textbf{D5:} Cấu hình Discord của lớp.
\item \textbf{D6:} Kết quả theo QĐ 1.20.
\end{itemize}
\textbf{Thuật toán:}
\begin{itemize}[leftmargin=*]
\item \textbf{B1:} Giáo viên chọn lớp và guild/channel.
\item \textbf{B2:} Kiểm tra lớp thuộc giáo viên.
\item \textbf{B3:} Kiểm tra guild/channel nằm trong dữ liệu Discord hệ thống nhìn thấy.
\item \textbf{B4:} Kiểm tra channel thuộc đúng guild.
\item \textbf{B5:} Lưu binding cho lớp.
\end{itemize}

\paragraph{Chức năng Gửi thông báo vào channel của lớp}
\dfddiagram{Gửi thông báo\\channel lớp}{Giáo viên}{Discord API}{ClassDiscordBinding, Log}
\textbf{Mô tả luồng dữ liệu:}
\begin{itemize}[leftmargin=*]
\item \textbf{D1:} Nội dung thông báo và danh sách guild nhận tin theo BM 1.21.
\item \textbf{D2:} Trạng thái gửi từ Discord.
\item \textbf{D3:} Binding lớp, bot credential và channel thông báo.
\item \textbf{D4:} Nhật ký gửi tin và trạng thái lỗi nếu có.
\item \textbf{D5:} Trạng thái gửi theo từng guild/channel.
\item \textbf{D6:} Kết quả theo QĐ 1.21.
\end{itemize}
\textbf{Thuật toán:}
\begin{itemize}[leftmargin=*]
\item \textbf{B1:} Nhận nội dung và danh sách guild.
\item \textbf{B2:} Kiểm tra nội dung không rỗng và có ít nhất một guild.
\item \textbf{B3:} Loại bỏ guild trùng.
\item \textbf{B4:} Gửi tin qua Discord API.
\item \textbf{B5:} Ghi nhận kết quả gửi.
\end{itemize}

\paragraph{Chức năng Gửi tin nhắn riêng cho học sinh qua Discord}
\dfddiagram{Gửi tin nhắn\\riêng Discord}{Giáo viên}{Discord API}{StudentDiscordCredential, Log}
\textbf{Mô tả luồng dữ liệu:}
\begin{itemize}[leftmargin=*]
\item \textbf{D1:} Nội dung, học sinh nhận tin hoặc lớp nhận tin theo BM 1.22.
\item \textbf{D2:} Trạng thái gửi DM từ Discord.
\item \textbf{D3:} Danh sách học sinh, StudentDiscordCredential và lớp học.
\item \textbf{D4:} Nhật ký gửi tin riêng.
\item \textbf{D5:} Kết quả gửi theo từng học sinh.
\item \textbf{D6:} Kết quả theo QĐ 1.22.
\end{itemize}
\textbf{Thuật toán:}
\begin{itemize}[leftmargin=*]
\item \textbf{B1:} Nhận nội dung và danh sách người nhận.
\item \textbf{B2:} Kiểm tra có ít nhất một học sinh hoặc một lớp.
\item \textbf{B3:} Lọc học sinh thuộc giáo viên hiện tại và có Discord credential.
\item \textbf{B4:} Gửi DM qua Discord API.
\item \textbf{B5:} Trả trạng thái gửi chi tiết.
\end{itemize}

\subsubsection{Quản lý Codeforces Gym}

\paragraph{Chức năng Gắn hoặc gỡ Gym Codeforces khỏi lớp học}
\dfddiagram{Gắn/gỡ Gym\\Codeforces}{Giáo viên}{Codeforces API}{Gym, GymProblem, ClassGym}
\textbf{Mô tả luồng dữ liệu:}
\begin{itemize}[leftmargin=*]
\item \textbf{D1:} Lớp học, Gym ID, chu kỳ đồng bộ theo BM 1.23.
\item \textbf{D2:} Dữ liệu Gym từ Codeforces; nếu chỉ gắn/gỡ Gym đã đồng bộ thì không có.
\item \textbf{D3:} Lớp, credential Codeforces, Gym hiện có.
\item \textbf{D4:} Liên kết Gym với lớp hoặc trạng thái gỡ Gym.
\item \textbf{D5:} Danh sách Gym của lớp.
\item \textbf{D6:} Kết quả theo QĐ 1.23.
\end{itemize}
\textbf{Thuật toán:}
\begin{itemize}[leftmargin=*]
\item \textbf{B1:} Giáo viên chọn lớp và Gym.
\item \textbf{B2:} Kiểm tra lớp thuộc giáo viên.
\item \textbf{B3:} Kiểm tra Gym ID là số và chu kỳ đồng bộ là số nguyên dương nếu có.
\item \textbf{B4:} Gắn/gỡ Gym, giữ dữ liệu lịch sử đã đồng bộ nếu cần báo cáo.
\end{itemize}

\paragraph{Chức năng Xem bảng xếp hạng Gym và hồ sơ học tập}
\dfddiagram{Xem bảng xếp hạng\\Gym}{Giáo viên}{Thiết bị nhập}{GymStanding, GymProblem, Student}
\textbf{Mô tả luồng dữ liệu:}
\begin{itemize}[leftmargin=*]
\item \textbf{D1:} Lớp học, Gym, học sinh hoặc bài cần xem.
\item \textbf{D2:} Không có.
\item \textbf{D3:} GymProblem, GymStanding, Student, Enrollment.
\item \textbf{D4:} Không có.
\item \textbf{D5:} Bảng xếp hạng Gym và hồ sơ học tập.
\item \textbf{D6:} Kết quả hiển thị hoặc thông báo chưa có dữ liệu.
\end{itemize}
\textbf{Thuật toán:}
\begin{itemize}[leftmargin=*]
\item \textbf{B1:} Nhận lớp/Gym cần xem.
\item \textbf{B2:} Kiểm tra lớp thuộc giáo viên.
\item \textbf{B3:} Đọc GymStanding đã đồng bộ.
\item \textbf{B4:} Ghép standing với học sinh đang ghi danh.
\item \textbf{B5:} Hiển thị bảng kết quả.
\end{itemize}

\subsubsection{Uỷ quyền Discord của học sinh}

\paragraph{Chức năng Học sinh uỷ quyền Discord OAuth}
\dfddiagram{Uỷ quyền Discord\\của học sinh}{Học sinh}{Discord OAuth}{StudentDiscordCredential}
\textbf{Mô tả luồng dữ liệu:}
\begin{itemize}[leftmargin=*]
\item \textbf{D1:} Liên kết uỷ quyền gắn với học sinh theo BM 1.24.
\item \textbf{D2:} OAuth consent, callback, Discord user id và token.
\item \textbf{D3:} Hồ sơ học sinh và credential hiện có.
\item \textbf{D4:} StudentDiscordCredential được lưu bảo mật.
\item \textbf{D5:} Trang thông báo uỷ quyền thành công/thất bại.
\item \textbf{D6:} Kết quả theo QĐ 1.24.
\end{itemize}
\textbf{Thuật toán:}
\begin{itemize}[leftmargin=*]
\item \textbf{B1:} Học sinh mở liên kết OAuth.
\item \textbf{B2:} Discord xác nhận quyền và trả callback.
\item \textbf{B3:} Hệ thống kiểm tra URL gắn với học sinh hợp lệ.
\item \textbf{B4:} Lưu Discord identity và token bảo mật.
\item \textbf{B5:} Trả trang kết quả.
\end{itemize}

\subsubsection{Quản trị hệ thống và tổng quan}

\paragraph{Chức năng Quản lý tài khoản giáo viên}
\dfddiagram{Quản lý tài khoản\\giáo viên}{Quản trị viên}{Thiết bị nhập}{Teacher}
\textbf{Mô tả luồng dữ liệu:}
\begin{itemize}[leftmargin=*]
\item \textbf{D1:} Giáo viên, tên đăng nhập, mật khẩu, trạng thái, Codeforces credential theo BM 1.25.
\item \textbf{D2:} Không có.
\item \textbf{D3:} Danh sách giáo viên và trạng thái hiện có.
\item \textbf{D4:} Tài khoản giáo viên được cập nhật hoặc vô hiệu hoá.
\item \textbf{D5:} Danh sách tài khoản giáo viên.
\item \textbf{D6:} Kết quả theo QĐ 1.25.
\end{itemize}
\textbf{Thuật toán:}
\begin{itemize}[leftmargin=*]
\item \textbf{B1:} Quản trị viên chọn tài khoản giáo viên.
\item \textbf{B2:} Kiểm tra quyền admin.
\item \textbf{B3:} Kiểm tra có ít nhất một trường thay đổi.
\item \textbf{B4:} Cập nhật hoặc vô hiệu hóa, không xóa cứng tài khoản.
\end{itemize}

\paragraph{Chức năng Xem tổng quan giảng dạy}
\dfddiagram{Xem tổng quan\\giảng dạy}{Giáo viên}{Thiết bị nhập}{Class, Session, Student, Finance}
\textbf{Mô tả luồng dữ liệu:}
\begin{itemize}[leftmargin=*]
\item \textbf{D1:} Yêu cầu mở dashboard hoặc bộ lọc thời gian.
\item \textbf{D2:} Không có.
\item \textbf{D3:} Lớp, học sinh, buổi học, điểm danh, tài chính, GymStanding.
\item \textbf{D4:} Không có.
\item \textbf{D5:} Số liệu tổng quan giảng dạy.
\item \textbf{D6:} Dashboard hoặc thông báo không có dữ liệu.
\end{itemize}
\textbf{Thuật toán:}
\begin{itemize}[leftmargin=*]
\item \textbf{B1:} Giáo viên mở dashboard.
\item \textbf{B2:} Hệ thống đọc dữ liệu thuộc giáo viên hiện tại.
\item \textbf{B3:} Tổng hợp số lớp, học sinh, buổi học, doanh thu, công nợ và kết quả Gym.
\item \textbf{B4:} Trả dữ liệu tổng quan.
\end{itemize}

\subsection{Sơ đồ sử dụng chức năng (Use Case Diagram)}

\subsubsection{Danh sách các Actor}

\begin{table}[!htbp]
\centering
\tablesetup
\begin{tabularx}{\textwidth}{|C{0.8cm}|L{2.6cm}|Y|}
\hline
\textbf{STT} & \textbf{Actor} & \textbf{Trách nhiệm} \\
\hline
1 & Giáo viên & Sử dụng hầu hết nghiệp vụ quản lý lớp học: học sinh, ghi danh, lớp, lịch, buổi học, điểm danh, tài chính, báo cáo, Discord và Codeforces Gym. \\
\hline
2 & Quản trị viên & Quản lý tài khoản giáo viên và cấu hình bot Discord toàn hệ thống. \\
\hline
3 & Học sinh & Uỷ quyền Discord OAuth để hệ thống có thể lưu danh tính Discord và phục vụ membership lớp học. \\
\hline
4 & Discord & Cung cấp OAuth, guild/channel/member API, voice state và gửi thông báo/DM. \\
\hline
5 & Codeforces & Cung cấp Gym, GymProblem và GymStanding cho quá trình theo dõi học tập. \\
\hline
\end{tabularx}
\end{table}

\subsubsection{Sơ đồ Use Case tổng quan}

\begin{figure}[!htbp]
\centering
\IfFileExists{images/use-case-overview.png}{%
    \includegraphics[width=0.96\textwidth]{images/use-case-overview.png}%
}{%
\IfFileExists{report/functional-requirements/diagrams/use-case-overview.png}{%
    \includegraphics[width=0.96\textwidth]{report/functional-requirements/diagrams/use-case-overview.png}%
}{%
    \fbox{\parbox[c][5cm][c]{0.9\textwidth}{\centering
    Chưa tìm thấy ảnh sơ đồ Use Case tổng quan.\\[4pt]
    Đặt file tại \code{images/use-case-overview.png} để Overleaf tự nạp.}}
}}%
\caption{Sơ đồ Use Case tổng quan của hệ thống giúp giáo viên quản lý lớp học lập trình thi đấu}
\end{figure}

\subsubsection{Đặc tả các Use Case}

\paragraph{Quản lý tài khoản}

\subparagraph{Đăng nhập hệ thống}
\begin{itemize}[leftmargin=*]
\item \textbf{Tóm tắt:} Cho phép giáo viên hoặc quản trị viên truy cập hệ thống theo đúng vai trò.
\item \textbf{Dòng sự kiện chính:} Người dùng nhập tên đăng nhập và mật khẩu; hệ thống kiểm tra tài khoản, mật khẩu và trạng thái; hệ thống cấp token và điều hướng vào màn hình phù hợp.
\item \textbf{Dòng sự kiện phụ:} Nếu thiếu thông tin, sai mật khẩu hoặc tài khoản bị vô hiệu hóa, hệ thống từ chối đăng nhập và hiển thị lỗi.
\item \textbf{Tiền điều kiện:} Tài khoản đã tồn tại và hệ thống đang hoạt động.
\end{itemize}

\subparagraph{Đăng ký tài khoản giáo viên}
\begin{itemize}[leftmargin=*]
\item \textbf{Tóm tắt:} Tạo tài khoản giáo viên mới để sử dụng hệ thống.
\item \textbf{Dòng sự kiện chính:} Người dùng nhập thông tin; hệ thống kiểm tra tên đăng nhập; hệ thống băm mật khẩu, lưu credential Codeforces nếu có và tạo tài khoản.
\item \textbf{Dòng sự kiện phụ:} Nếu tên đăng nhập trùng hoặc thiếu dữ liệu bắt buộc, hệ thống báo lỗi.
\item \textbf{Tiền điều kiện:} Chức năng đăng ký hoặc quyền quản trị tạo tài khoản đang được cho phép.
\end{itemize}

\subparagraph{Cập nhật hồ sơ giáo viên}
\begin{itemize}[leftmargin=*]
\item \textbf{Tóm tắt:} Cho phép giáo viên hoặc quản trị viên cập nhật tên đăng nhập, mật khẩu và thông tin Codeforces.
\item \textbf{Dòng sự kiện chính:} Người dùng nhập trường cần đổi; hệ thống kiểm tra ràng buộc; hệ thống lưu hồ sơ và che dữ liệu nhạy cảm.
\item \textbf{Dòng sự kiện phụ:} Nếu không có trường thay đổi hoặc mật khẩu rỗng, hệ thống báo lỗi.
\item \textbf{Tiền điều kiện:} Người dùng đã đăng nhập hoặc quản trị viên đang thao tác trên tài khoản giáo viên.
\end{itemize}

\subparagraph{Liên kết Discord của giáo viên}
\begin{itemize}[leftmargin=*]
\item \textbf{Tóm tắt:} Cho phép giáo viên xác minh Discord để dùng các chức năng guild/channel.
\item \textbf{Dòng sự kiện chính:} Giáo viên yêu cầu liên kết; hệ thống sinh URL OAuth; Discord trả callback; hệ thống lưu Discord user id và username.
\item \textbf{Dòng sự kiện phụ:} Nếu callback không hợp lệ hoặc Discord từ chối quyền, hệ thống không lưu dữ liệu.
\item \textbf{Tiền điều kiện:} Giáo viên đã đăng nhập và cấu hình OAuth Discord hợp lệ.
\end{itemize}

\paragraph{Quản lý học sinh và ghi danh}

\subparagraph{Thêm học sinh và ghi danh vào lớp}
\begin{itemize}[leftmargin=*]
\item \textbf{Tóm tắt:} Tạo hồ sơ học sinh và ghi danh học sinh vào một lớp của giáo viên.
\item \textbf{Dòng sự kiện chính:} Giáo viên nhập thông tin; hệ thống kiểm tra lớp thuộc giáo viên, dữ liệu bắt buộc và ghi danh trùng; hệ thống tạo Student và Enrollment.
\item \textbf{Dòng sự kiện phụ:} Nếu lớp không thuộc giáo viên hoặc thiếu Codeforces handle, hệ thống báo lỗi.
\item \textbf{Tiền điều kiện:} Giáo viên đã đăng nhập và lớp học đã tồn tại.
\end{itemize}

\subparagraph{Cập nhật hồ sơ học sinh}
\begin{itemize}[leftmargin=*]
\item \textbf{Tóm tắt:} Cập nhật thông tin cá nhân và Codeforces handle của học sinh.
\item \textbf{Dòng sự kiện chính:} Giáo viên chọn học sinh; nhập thông tin cần đổi; hệ thống kiểm tra quyền sở hữu; hệ thống lưu hồ sơ.
\item \textbf{Dòng sự kiện phụ:} Nếu học sinh không thuộc giáo viên hoặc không có trường thay đổi, hệ thống báo lỗi.
\item \textbf{Tiền điều kiện:} Học sinh đã tồn tại trong hệ thống.
\end{itemize}

\subparagraph{Chuyển lớp, rút học sinh, khôi phục học sinh}
\begin{itemize}[leftmargin=*]
\item \textbf{Tóm tắt:} Điều chỉnh trạng thái ghi danh và vòng đời học sinh.
\item \textbf{Dòng sự kiện chính:} Giáo viên chọn học sinh và thao tác; hệ thống kiểm tra lớp/học sinh thuộc giáo viên; hệ thống cập nhật Enrollment và trạng thái học sinh.
\item \textbf{Dòng sự kiện phụ:} Nếu học sinh còn số dư khi rút, hệ thống chuyển sang trạng thái phù hợp để xử lý công nợ.
\item \textbf{Tiền điều kiện:} Học sinh có ghi danh hoặc trạng thái có thể khôi phục.
\end{itemize}

\subparagraph{Lưu trữ học sinh đang chờ xử lý công nợ}
\begin{itemize}[leftmargin=*]
\item \textbf{Tóm tắt:} Lưu trữ học sinh sau khi đã xử lý xong tình trạng công nợ.
\item \textbf{Dòng sự kiện chính:} Giáo viên chọn học sinh chờ lưu trữ; hệ thống kiểm tra số dư; hệ thống chuyển sang trạng thái đã lưu trữ.
\item \textbf{Dòng sự kiện phụ:} Nếu còn số dư cần thu hoặc hoàn, hệ thống hiển thị lý do chưa thể hoàn tất.
\item \textbf{Tiền điều kiện:} Học sinh đang ở trạng thái chờ lưu trữ.
\end{itemize}

\subparagraph{Tra cứu và xem hồ sơ học tập}
\begin{itemize}[leftmargin=*]
\item \textbf{Tóm tắt:} Tra cứu học sinh và xem dữ liệu học tập, tài chính, Codeforces Gym.
\item \textbf{Dòng sự kiện chính:} Giáo viên nhập bộ lọc; hệ thống tìm học sinh thuộc giáo viên; hệ thống tổng hợp hồ sơ, ghi danh, điểm danh, công nợ và standing Gym.
\item \textbf{Dòng sự kiện phụ:} Nếu không có dữ liệu phù hợp, hệ thống hiển thị thông báo rỗng.
\item \textbf{Tiền điều kiện:} Đã có dữ liệu học sinh.
\end{itemize}

\paragraph{Quản lý lớp học, lịch học và buổi học}

\subparagraph{Thêm/sửa/lưu trữ lớp học}
\begin{itemize}[leftmargin=*]
\item \textbf{Tóm tắt:} Quản lý lớp học và trạng thái lớp.
\item \textbf{Dòng sự kiện chính:} Giáo viên nhập tên lớp, học phí và lịch học; hệ thống kiểm tra dữ liệu; hệ thống lưu lớp hoặc chuyển trạng thái archived.
\item \textbf{Dòng sự kiện phụ:} Nếu học phí âm hoặc thiếu lịch khi tạo lớp, hệ thống báo lỗi.
\item \textbf{Tiền điều kiện:} Giáo viên đã đăng nhập.
\end{itemize}

\subparagraph{Thiết lập lịch học định kỳ}
\begin{itemize}[leftmargin=*]
\item \textbf{Tóm tắt:} Cấu hình các khung giờ học định kỳ cho lớp.
\item \textbf{Dòng sự kiện chính:} Giáo viên nhập danh sách lịch; hệ thống kiểm tra thứ và giờ; hệ thống thay thế lịch đang áp dụng.
\item \textbf{Dòng sự kiện phụ:} Nếu giờ kết thúc không sau giờ bắt đầu, hệ thống báo lỗi.
\item \textbf{Tiền điều kiện:} Lớp học thuộc giáo viên hiện tại.
\end{itemize}

\subparagraph{Tạo buổi học thủ công và xem danh sách buổi học}
\begin{itemize}[leftmargin=*]
\item \textbf{Tóm tắt:} Tạo buổi học phát sinh hoặc xem các buổi học đã có.
\item \textbf{Dòng sự kiện chính:} Giáo viên chọn lớp, nhập thời gian; hệ thống kiểm tra thời gian; hệ thống tạo Session hoặc trả danh sách Session.
\item \textbf{Dòng sự kiện phụ:} Nếu thiếu nhóm dữ liệu thời gian hợp lệ, hệ thống báo lỗi.
\item \textbf{Tiền điều kiện:} Lớp học tồn tại và đang thuộc giáo viên.
\end{itemize}

\subparagraph{Hủy buổi học}
\begin{itemize}[leftmargin=*]
\item \textbf{Tóm tắt:} Chuyển buổi học sang trạng thái hủy mà không xóa lịch sử.
\item \textbf{Dòng sự kiện chính:} Giáo viên chọn buổi học; hệ thống kiểm tra quyền; hệ thống cập nhật trạng thái hủy và xử lý dữ liệu liên quan.
\item \textbf{Dòng sự kiện phụ:} Nếu buổi học không thuộc giáo viên, hệ thống từ chối.
\item \textbf{Tiền điều kiện:} Buổi học đã tồn tại.
\end{itemize}

\paragraph{Quản lý điểm danh}

\subparagraph{Xem và cập nhật trạng thái điểm danh học sinh}
\begin{itemize}[leftmargin=*]
\item \textbf{Tóm tắt:} Ghi nhận có mặt, vắng có phép hoặc vắng không phép cho từng học sinh.
\item \textbf{Dòng sự kiện chính:} Giáo viên mở buổi học; hệ thống tải danh sách học sinh; giáo viên cập nhật trạng thái; hệ thống lưu Attendance và cập nhật FeeRecord.
\item \textbf{Dòng sự kiện phụ:} Nếu buổi học đã hủy hoặc trạng thái không hợp lệ, hệ thống từ chối lưu.
\item \textbf{Tiền điều kiện:} Buổi học và học sinh thuộc giáo viên hiện tại.
\end{itemize}

\paragraph{Quản lý tài chính và báo cáo}

\subparagraph{Ghi nhận giao dịch thu tiền hoặc hoàn tiền}
\begin{itemize}[leftmargin=*]
\item \textbf{Tóm tắt:} Lưu giao dịch thu học phí hoặc hoàn tiền cho học sinh.
\item \textbf{Dòng sự kiện chính:} Giáo viên nhập giao dịch; hệ thống kiểm tra học sinh và loại giao dịch; hệ thống lưu Transaction và cập nhật số dư.
\item \textbf{Dòng sự kiện phụ:} Nếu học sinh không thuộc giáo viên hoặc loại giao dịch không hợp lệ, hệ thống báo lỗi.
\item \textbf{Tiền điều kiện:} Học sinh đã tồn tại.
\end{itemize}

\subparagraph{Cập nhật trạng thái khoản phí}
\begin{itemize}[leftmargin=*]
\item \textbf{Tóm tắt:} Đánh dấu khoản phí đang hiệu lực hoặc đã hủy.
\item \textbf{Dòng sự kiện chính:} Giáo viên chọn khoản phí; hệ thống kiểm tra quyền; hệ thống cập nhật trạng thái FeeRecord.
\item \textbf{Dòng sự kiện phụ:} Nếu trạng thái không hợp lệ, hệ thống báo lỗi.
\item \textbf{Tiền điều kiện:} Khoản phí đã phát sinh từ điểm danh hoặc buổi học.
\end{itemize}

\subparagraph{Xem số dư và tổng quan tài chính}
\begin{itemize}[leftmargin=*]
\item \textbf{Tóm tắt:} Cho phép giáo viên đối soát công nợ theo học sinh hoặc lớp.
\item \textbf{Dòng sự kiện chính:} Giáo viên chọn bộ lọc; hệ thống tổng hợp FeeRecord và Transaction; hệ thống hiển thị số dư.
\item \textbf{Dòng sự kiện phụ:} Nếu không có dữ liệu, hệ thống hiển thị bảng rỗng.
\item \textbf{Tiền điều kiện:} Có dữ liệu học sinh, khoản phí hoặc giao dịch.
\end{itemize}

\subparagraph{Lập báo cáo thu nhập theo kỳ}
\begin{itemize}[leftmargin=*]
\item \textbf{Tóm tắt:} Tổng hợp thu nhập, khoản phải thu và số dư theo khoảng thời gian.
\item \textbf{Dòng sự kiện chính:} Giáo viên nhập bộ lọc; hệ thống kiểm tra thời gian và lớp; hệ thống tổng hợp báo cáo.
\item \textbf{Dòng sự kiện phụ:} Nếu từ ngày sau đến ngày, hệ thống yêu cầu nhập lại.
\item \textbf{Tiền điều kiện:} Giáo viên đã đăng nhập.
\end{itemize}

\paragraph{Quản lý Discord}

\subparagraph{Cấu hình bot Discord toàn hệ thống}
\begin{itemize}[leftmargin=*]
\item \textbf{Tóm tắt:} Quản trị viên lưu credential bot và OAuth cho hệ thống.
\item \textbf{Dòng sự kiện chính:} Quản trị viên nhập credential; hệ thống kiểm tra quyền admin; hệ thống lưu bảo mật.
\item \textbf{Dòng sự kiện phụ:} Nếu thiếu token hoặc client secret, hệ thống báo lỗi.
\item \textbf{Tiền điều kiện:} Người dùng có vai trò quản trị viên.
\end{itemize}

\subparagraph{Chọn Discord guild và channel cho lớp}
\begin{itemize}[leftmargin=*]
\item \textbf{Tóm tắt:} Gắn guild, channel thông báo và voice channel điểm danh vào lớp.
\item \textbf{Dòng sự kiện chính:} Giáo viên chọn lớp, guild và channel; hệ thống kiểm tra dữ liệu Discord; hệ thống lưu ClassDiscordBinding.
\item \textbf{Dòng sự kiện phụ:} Nếu channel không thuộc guild đã chọn, hệ thống báo lỗi.
\item \textbf{Tiền điều kiện:} Giáo viên đã liên kết Discord và bot đã được cấu hình.
\end{itemize}

\subparagraph{Gửi thông báo vào channel của lớp}
\begin{itemize}[leftmargin=*]
\item \textbf{Tóm tắt:} Gửi thông báo chung tới channel Discord đã cấu hình.
\item \textbf{Dòng sự kiện chính:} Giáo viên nhập nội dung; hệ thống kiểm tra guild/channel; hệ thống gửi qua Discord API và ghi trạng thái.
\item \textbf{Dòng sự kiện phụ:} Nếu Discord lỗi hoặc chưa có channel, hệ thống trả trạng thái lỗi.
\item \textbf{Tiền điều kiện:} Lớp đã có cấu hình Discord hợp lệ.
\end{itemize}

\subparagraph{Gửi tin nhắn riêng cho học sinh qua Discord}
\begin{itemize}[leftmargin=*]
\item \textbf{Tóm tắt:} Gửi DM tới học sinh đã uỷ quyền Discord.
\item \textbf{Dòng sự kiện chính:} Giáo viên nhập nội dung và chọn người nhận; hệ thống lọc học sinh có credential; hệ thống gửi DM và trả trạng thái từng học sinh.
\item \textbf{Dòng sự kiện phụ:} Nếu học sinh chưa có Discord credential, hệ thống bỏ qua hoặc báo lý do.
\item \textbf{Tiền điều kiện:} Học sinh đã uỷ quyền Discord.
\end{itemize}

\paragraph{Quản lý Codeforces Gym}

\subparagraph{Gắn hoặc gỡ Gym Codeforces khỏi lớp học}
\begin{itemize}[leftmargin=*]
\item \textbf{Tóm tắt:} Liên kết Gym với lớp để theo dõi bài và standing.
\item \textbf{Dòng sự kiện chính:} Giáo viên chọn lớp và Gym; hệ thống kiểm tra quyền; hệ thống lưu liên kết hoặc gỡ liên kết.
\item \textbf{Dòng sự kiện phụ:} Nếu Gym ID không phải số hoặc credential lỗi, hệ thống báo lỗi.
\item \textbf{Tiền điều kiện:} Lớp học tồn tại và giáo viên có credential Codeforces phù hợp khi cần đồng bộ.
\end{itemize}

\subparagraph{Xem bảng xếp hạng Gym và hồ sơ học tập}
\begin{itemize}[leftmargin=*]
\item \textbf{Tóm tắt:} Hiển thị standing theo học sinh đang học trong lớp.
\item \textbf{Dòng sự kiện chính:} Giáo viên chọn lớp/Gym; hệ thống đọc GymStanding; hệ thống ghép với danh sách học sinh và hiển thị bảng kết quả.
\item \textbf{Dòng sự kiện phụ:} Nếu Gym chưa đồng bộ, hệ thống thông báo chưa có dữ liệu.
\item \textbf{Tiền điều kiện:} Lớp đã gắn Gym.
\end{itemize}

\paragraph{Uỷ quyền Discord của học sinh}

\subparagraph{Học sinh uỷ quyền Discord OAuth}
\begin{itemize}[leftmargin=*]
\item \textbf{Tóm tắt:} Học sinh cấp quyền để hệ thống lưu Discord identity và token.
\item \textbf{Dòng sự kiện chính:} Học sinh mở URL; Discord xác nhận quyền; hệ thống nhận callback, kiểm tra học sinh và lưu credential.
\item \textbf{Dòng sự kiện phụ:} Nếu callback sai hoặc học sinh không tồn tại, hệ thống không lưu credential.
\item \textbf{Tiền điều kiện:} Học sinh có liên kết uỷ quyền hợp lệ.
\end{itemize}

\paragraph{Quản trị hệ thống}

\subparagraph{Quản lý tài khoản giáo viên}
\begin{itemize}[leftmargin=*]
\item \textbf{Tóm tắt:} Quản trị viên cập nhật hoặc vô hiệu hóa tài khoản giáo viên.
\item \textbf{Dòng sự kiện chính:} Quản trị viên chọn tài khoản; nhập trường cần thay đổi; hệ thống kiểm tra quyền; hệ thống cập nhật tài khoản.
\item \textbf{Dòng sự kiện phụ:} Nếu không có trường thay đổi, hệ thống báo lỗi.
\item \textbf{Tiền điều kiện:} Người dùng đang đăng nhập với vai trò quản trị viên.
\end{itemize}

\subparagraph{Xem tổng quan giảng dạy}
\begin{itemize}[leftmargin=*]
\item \textbf{Tóm tắt:} Giáo viên xem nhanh số liệu lớp, học sinh, buổi học, tài chính và kết quả học tập.
\item \textbf{Dòng sự kiện chính:} Giáo viên mở dashboard; hệ thống tổng hợp dữ liệu thuộc giáo viên; hệ thống hiển thị chỉ số tổng quan.
\item \textbf{Dòng sự kiện phụ:} Nếu chưa có dữ liệu, hệ thống hiển thị trạng thái rỗng.
\item \textbf{Tiền điều kiện:} Giáo viên đã đăng nhập.
\end{itemize}

\end{document}
